\documentclass[11pt]{amsart}
\usepackage[T1]{fontenc}
\usepackage{lmodern}
\usepackage{amsmath,amssymb,mathtools,mathrsfs,booktabs,longtable,microtype}
\usepackage[margin=1.08in]{geometry}
\usepackage[hidelinks]{hyperref}
\hypersetup{pdftitle={General Theta Foundations I: Acquired Geometry and Causal Resource Transfer},pdfauthor={Qian Qi}}
\pdfinfoomitdate=1
\pdftrailerid{}
\pdfsuppressptexinfo=15
\newtheorem{theorem}{Theorem}[section]
\newtheorem{proposition}[theorem]{Proposition}
\newtheorem{lemma}[theorem]{Lemma}
\newtheorem{corollary}[theorem]{Corollary}
\theoremstyle{definition}
\newtheorem{definition}[theorem]{Definition}
\newtheorem{example}[theorem]{Example}
\theoremstyle{remark}
\newtheorem{remark}[theorem]{Remark}
\newcommand{\E}{\mathbb E}
\newcommand{\PP}{\mathbb P}
\newcommand{\R}{\mathbb R}
\newcommand{\N}{\mathbb N}
\newcommand{\TV}{\operatorname{TV}}
\newcommand{\law}{\operatorname{Law}}
\newcommand{\osc}{\operatorname{osc}}
\newcommand{\dist}{\operatorname{dist}}
\newcommand{\cp}{\mathrm{cp}}
\newcommand{\on}{\mathrm{on}}
\newcommand{\one}{\mathbf 1}
\newcommand{\norm}[1]{\left\lVert#1\right\rVert}
\newcommand{\abs}[1]{\left\lvert#1\right\rvert}
\title{General Theta Foundations I:\\Acquired Geometry and Causal Resource Transfer}
\author{Qian Qi}
\date{7 October 2026}
\subjclass[2020]{Primary 60G35, 62B05; Secondary 94A34, 93E11, 28A80}
\keywords{causal experiment, predictive quotient, acquired geometry, continuation information, adaptive acquisition, finite memory}
\begin{document}
\begin{abstract}
A prepared causal experiment acquires a geometry of conditional future responses. We prove a two-sided finite-resource theorem for experiments whose observable acquired measures admit global recursive covers. Their optimal online excess is governed by the largest continuation exponent across all information cuts; checkpoint excess is governed by the terminal exponent alone. This gives a necessary and sufficient criterion for their constant-factor agreement within the stated certificate class. The proof combines actual submeasure lower bounds with a candidate-moment recursion, allowing dependent reports, changing rank and singular acquisition. A positive-noise regression experiment has online exponent determined by observable rank rather than hidden dimension. A singular observation experiment exhibits a transition when a continuous future query changes the terminal geometry. We also establish a matched acquisition--memory--calibration law with an attained causal defect for progressive sensors. Numerical error and resource-certified experiment morphisms compose in the same scored task. Global characterization of arbitrary causal encoders is not assumed or asserted.
\end{abstract}
\maketitle
\section{The experimental problem}\label{sec:problem}
A causal experiment does not present an observer with a parameter manifold. It presents preparations, admissible interventions and reports. Geometry becomes relevant only after specifying which future reports are to be predicted and which information the observer has actually acquired. There is a further obstruction: a small terminal codebook need not be recursively selectable after earlier reports have been discarded. We use \emph{General Theta} for the mathematical program connecting prepared causal kernels, executable future tests, their minimal predictive quotient, resource-certified experiment comparison, acquired geometry and prediction or decision risk. ``Foundations'' refers to this order of construction, not to a claim that all causal encoders have already been classified.

The central result is Theorem~\ref{thm:serial}. Suppose the acquired response at cut $t$ has a positive-mass observable submeasure of exponent $s_t$, and a globally selectable cover of the same order is stable under the raw task updates. The stability requirement includes a Lyapunov envelope for all candidates, not just the exact trajectory or the code eventually chosen. Then, for one fixed terminal task and its full-history Bayes risk $B_T$,
\[
 R_\on(T,M)-B_T\asymp M^{-2/s_*},\qquad
 R_\cp(T,M)-B_T\asymp M^{-2/s_T},\qquad s_*=\max_t s_t.
\]
The entire family of actual continuation widths is equivalent in order to the online risk. Within this declared class, checkpoint and online compression match if and only if $s_*=s_T$. The theorem allows noninteger exponents, changing state dimension, dependent reports and expansive finite-horizon updates. Its proof constructs one $M$-label recursion at every cut; it does not assume that separately optimal checkpoint codebooks can be combined.

Two mechanisms make this possible. A conditional future response is a function of the information arriving after a cut. An earlier label must therefore choose a function, not merely a terminal scalar. Quantizing the actual dominated measure of these functions yields a common-task lower. Theorem~\ref{thm:cutduality} identifies the exact one-cut functional-coding problem, compares it with acquired Hilbert measures, and localizes it when physically shared side information prevents global product domination. Raw density and joint regeneration lemmas produce its common continuation laws. On the constructive side, global covers and a candidate Lyapunov bound control every inserted error under the actual report law, without independence or a free exact posterior register. The block transfer theorem, Theorem~\ref{thm:transfer}, retains its separate strength of time-uniform control under verified raw-block stability.

Theorem~\ref{thm:ranknoisy} verifies the central theorem from positive-noise correlated Gaussian kernels. Its successive observable ranks $r_j$ can change, and its largest rank $r$ is determined by the executable query family. It proves online excess $M^{-2/r}$ against checkpoint excess $M^{-2}$, even in an ambient hidden dimension $d>r$. A partial task statistic is not mislabeled a full predictive quotient: nuisance directions can still affect the suffix law. Theorem~\ref{thm:singularserial} instead starts from a Cantor-product preparation of total exponent $D$, obtained from its actual cylinder masses. A noisy continuation and continuous terminal query produce online excess $M^{-2/\max(D,1)}$. The terminal cut is essential when $D<1$. These two realizations are consequences of the same general theorem, not premises from which it is inferred.

Acquisition itself is a different optimization. A fixed exploration lower cannot be optimized by interchanging infima over its policy-dependent acquired law. Theorem~\ref{thm:adaptive} resolves this problem for progressive sensors with independent fair unread tails and arbitrary nonhomogeneous scales. The observed bit values do not alter future tail scales; this is why deterministic greedy allocation attains the optimum over all admitted adaptive allocations. Theorem~\ref{thm:joint} matches acquisition, retained memory, calibration and an attained simulation defect in one task. Theorem~\ref{thm:digital} gives a finite-bit feasible account, not a general computational lower. Theorem~\ref{thm:brier} preserves the sharp persistent-state law for ordinary filtering loss in sparse hidden-state models. A sequential-chart theorem and a full-rank noisy regression theorem are retained as intermediate consequences and independently checkable specializations.

The compatibility criterion requires both acquired observability and globally implementable update certificates. It does not derive the latter from bare one-cut widths for every causal experiment. The matched acquisition theorem does not solve unknown-kernel or arbitrary observation-dependent experimental design. Conditional projection, functional source coding, quantization, classical allocation and experiment comparison are explicitly credited in Section~\ref{sec:literature}. Supplement S~\cite{companion} is part of this same submission, supplied in full; it preserves the nonreset singular, stopped, calibration, output-grid and typed-interface proofs. It is not an external publication or a repository-only proof premise.

\subsection{Kernels, interfaces and resources}
All measurable spaces are standard Borel. An experiment is specified by a preparation $\mu_\lambda$, hidden states $X_t$, actions $A_t$, reports $Y_{t+1}$, and jointly measurable probability kernels
\begin{equation}\label{eq:kernel}
 K_t^\lambda(dx',dy,dc\mid x,a).
\end{equation}
Here $c$ records nonnegative raw-call, time and other declared resource increments. Positivity means nonnegativity and normalization. Dirac operations, failure symbols and missing reports are allowed. A visible history $h_t$ consists of the preparation report, public exploration coins, and intervening actions, reports and public costs. Private exploration coins that influence later actions remain in the physical/controller state and the actual conditional law unless genuinely reported. Legal action sets have a declared measurable graph. Strategies are parameter-independent Borel kernels on those sets; stopping rules use the visible filtration. They cannot consult future reports or the unknown $\lambda$.

A continuation interface is a measurable function $v_t(h_t)$ which specifies legal future actions, available resources and any publicly supplied phase. A predictive realization is a standard Borel state $S_t=s_t(h_t)$ determining that interface, every declared future-test expectation, and jointly Borel report and update kernels $J_t,F_t$. When a phase is kept separate, write the state space as a measurable disjoint union $S=\coprod_{v\in V}S_v$ over a \emph{finite} public phase set. All comparison pairs at a given time lie in the same fiber $S_v$. Updates map both members into the same next fiber, whose phase is a declared function of the old phase, action and common report. No unretained data-dependent phase is public for free. An unbounded or data-dependent external interface requires a separately stated conditional theorem and is not covered by the alphabet count below.

A predictor retains at most $M$ data labels at each declared cut. Its update uses only its current retained tuple, current action/report and allowed fresh randomness. Any controller, simulator or clock register consulted later is counted separately or included in the label tuple. Temporary workspace must be erased before the next cut; program constants must be known, parameter-independent data with a charged description. A vector output is a write-only stream: previously emitted coordinates are not rereadable workspace. The notation $M$ is an alphabet cardinality, not total arithmetic space. If a public finite phase has $Q$ states, the full serial cardinality is at most $MQ$; a lower bound at a cut uses the \emph{entire} accessible data-dependent tuple. Deterministic time is fixed in the main prediction statements and is not a hidden observation channel.

For the general transfer theorem fix a preparation and a legal exploration independently of the prediction encoder. Algorithm coins are independent of that experiment. In the adaptive acquisition theorem the exploration is instead optimized together with the encoder; its lower bounds are proved anew and do not condition an encoder-dependent law as though it were fixed.

\subsection{Executable tests and a measurable predictive quotient}
An executable test is a legal finite continuation followed by a bounded function of its scored reports. Two histories are equivalent if their continuation interfaces and all test expectations agree. A set-theoretic quotient need not be a suitable measurable state. We record a sufficient realization result.

\begin{lemma}[A compact measurable section]\label{lem:selector}
A continuous surjection $e:K\to S$ between compact metrizable spaces has a Borel section.
\end{lemma}
\begin{proof}
For each $n$ fix a finite closed cover $(F_{n,j})_j$ of $K$ by nonempty sets of diameter at most $2^{-n}$. Inductively select the first $j_n$ for which $e^{-1}(s)\cap F_{1,j_1}\cap\cdots\cap F_{n,j_n}$ is nonempty. For a fixed finite index word its compact intersection has closed image under $e$. Each finite-word selection event is therefore Borel, using finite Boolean combinations of those closed images. The selected compact intersections are nested, meet the fiber, and have diameter tending to zero. They have a unique common point in $e^{-1}(s)$. Choosing once a point of each nonempty finite-word intersection gives Borel finite-valued approximants converging to that point. Their pointwise limit is a Borel section.
\end{proof}

\begin{proposition}[Continuous-instrument realization]\label{prop:quotient}
Let $\mathcal B$ be compact metrizable and let $\mathcal V\subset C(\mathcal B)$ be the separable closed linear span of constants and the continuous evaluations of a determining family of executable tests and interface coordinates. All kernel, density, update and instrument versions below are jointly Borel also in the action variable on the declared standard Borel action--report domain. For each fixed action, let its report space be Polish, with reference measure $m_a$, and let $g_a(\pi,y)$ be a jointly continuous density. On the open set $D_a=\{(\pi,y):g_a(\pi,y)>0\}$, with the relative product topology, suppose the normalized update $B_a$ is continuous. For every $f\in\mathcal V$, require the instrument evaluation $g_a(\pi,y)f(B_a(\pi,y))$, with its specified value at density zero, to belong to $\mathcal V$ as a function of $\pi$ and to be jointly continuous. Then the countable evaluation image is a compact metrizable predictive quotient. Report laws and updates descend, and the latter are continuous on the descended open positive-density domain. Every standard Borel statistic measurably determining all these evaluations measurably determines the quotient.
\end{proposition}
\begin{proof}
Normalize a countable dense family of evaluations and map $\mathcal B$ into $[-1,1]^{\N}$. The image $S$ is compact, hence Borel. Equality of its coordinates gives equality on $\mathcal V$ by uniform approximation and hence on the determining tests. Conversely, equality of all future tests and interface coordinates gives equality on their closed linear span. Thus the fibers are exactly predictive equivalence classes, not a finer partition carrying unrelated coordinates. Taking $f=1$ shows that density values agree on each fiber. Dividing equal instrument values by that positive common density shows that the evaluation of the updated belief agrees on the fiber.

The evaluation map is a continuous surjection from a compact space to a Hausdorff space, hence is perfect. Its product with the report identity is also a quotient map: closedness follows by taking a convergent subnet in the compact belief domain whenever an image net converges. The instrument coordinates and their positive-density ratios therefore descend continuously. The positive-density domain is open in $S\times Y$; no update is defined by division at a null report. For measurability jointly across actions, Lemma~\ref{lem:selector} supplies a Borel section of the compact evaluation map. Composing the assumed joint versions with this section supplies joint Borel versions of all descended kernels. A common jointly Borel null-report extension, when needed on wrong candidates, is additional data.

For minimality, collect the measurable coordinate factors through the given statistic into a measurable product map. It takes values in the Borel set $S$ on the domain in question; extend outside that set by a fixed point. Countably many almost-sure factors can be imposed on a single common full-measure history set. This proves measurable, rather than merely set-theoretic, minimality.
\end{proof}

Applications may verify a predictive realization directly. The metric used for a quantitative estimate need only induce its standard Borel structure; it need not be the compact topology in Proposition~\ref{prop:quotient}.

\subsection{One prediction task}
Let $Y$ be a bounded terminal vector in a finite-dimensional Euclidean space with a specified weighted norm. A finite independent terminal probe index is interpreted as the vector of its separately executable conditional means, with the probe probabilities as weights; it is not a simultaneous counterfactual measurement. At the decision cut $T$ put
\begin{equation}\label{eq:task}
 P_T=\E[Y\mid\mathcal H_T],\qquad B_T=\E\norm{Y-P_T}^2,
 \qquad \nu_T=\law(P_T).
\end{equation}
For a finite Borel measure $\nu$ on a separable Hilbert space define
\begin{equation}\label{eq:quantization}
 q_M(\nu)=\inf_{c_1,\ldots,c_M}\int\min_i\norm{x-c_i}^2\,\nu(dx).
\end{equation}
The measure may be atomic, singular or a subprobability. A checkpoint algorithm inspects $\mathcal H_T$ once before retaining $M$ labels. An online algorithm retains at most $M$ labels at every earlier cut and cannot reread a report. In both cases the terminal outcome is revealed only after the prediction. Denote the optimal risks by $R_\cp(T,M)$ and $R_\on(T,M)$.

\section{Acquired continuation geometry and causal transfer}\label{sec:transfer}
We first specify the lower certificate, which depends only on the physical experiment and its fixed continuation, not on a proposed encoder.

\begin{definition}[Dominated continuation cut]\label{def:cut}
At a deterministic cut $s\le T$, let $H$ be the actual prefix and $U$ the complete physical pre-decision suffix, including all post-cut public actions, phases, costs and reports. The complete decision history is generated by $(H,U)$. The fixed exploration supplies their joint law $\Lambda(dh,du)$. A dominated continuation certificate consists of a probability measure $\eta$ on the suffix space and a finite measure $\underline\mu$ on prefix histories such that
\begin{equation}\label{eq:domination}
 \underline\mu(dh)\eta(du)\le\Lambda(dh,du).
\end{equation}
Choose a bounded jointly Borel version $z(h,u)=\E[Y\mid H=h,U=u]$ and set
\[
 Z_s(h)=z(h,\cdot)\in L^2(\eta;\R^q),\qquad
 \xi_s=(Z_s)_*\underline\mu.
\]
The task norm supplies the inner product in this function space. All post-cut information supplied to a decoder is contained in $U$ and independent algorithm coins; no earlier tape or additional data-dependent public register remains outside its counted label.
\end{definition}

The product inequality is a raw-law minorization, not independence of the actual suffix. For example a density $j(u\mid h)\ge c$ relative to $\eta$ on an actual prefix event $E$ gives $\underline\mu=c\,\law(H)|_E$. An independent suffix has $c=1$. In the absence of such a nonzero certificate the corresponding bound is zero, not an assertion that the experiment is easy. Domination ensures that versions of $z$ give the same element for $\underline\mu$-almost every $h$. Standard Borel spaces have countably generated sigma fields modulo a fixed probability, so $L^2(\eta)$ is separable. Approximation of a bounded jointly measurable function by simple rectangle functions, followed by dominated convergence, makes $h\mapsto Z_s(h)$ a Borel map into this Hilbert space.

\begin{lemma}[Projection and the continuation cut]\label{lem:cut}
For the fixed task in \eqref{eq:task},
\begin{equation}\label{eq:checkpoint}
 R_\cp(T,M)=B_T+q_M(\nu_T).
\end{equation}
Every $M$-label online predictor and every certificate in Definition~\ref{def:cut} satisfy
\begin{equation}\label{eq:cutlower}
 R_\on(T,M)-B_T\ge q_M(\xi_s).
\end{equation}
Both assertions allow independent shared and private randomization. They use the same terminal task and baseline.
\end{lemma}
\begin{proof}
Conditional projection gives $\E\norm{Y-A}^2=B_T+\E\norm{P_T-A}^2$. Fix all algorithm coins, which are independent of the fixed exploration. A checkpoint readout has at most $M$ centers. Randomized readouts may alternatively be replaced by conditional means, and nearest-center selection is Borel. Taking approximating center lists proves \eqref{eq:checkpoint}; averaging independent seeds does not change $\nu_T$.

At cut $s$ write its retained label as $i(H)$. Give the suffix decoder all of $U$, even when its original online implementation would retain less. Conditional on fixed coins, including all future private decoder coins, composing the finite sequence of later Borel updates gives an output $a_{i(H)}(U)$ for at most $M$ Borel response functions. Here $U$ contains all post-cut public actions, phases, costs and reports, not observation values alone. Projection and \eqref{eq:domination} give
\[
 R-B_T\ge\int\norm{z(h,u)-a_{i(h)}(u)}^2
                   \underline\mu(dh)\eta(du)
 \ge\int\min_i\norm{Z_s(h)-a_i}_{L^2(\eta)}^2\underline\mu(dh).
\]
Outputs can be projected into a fixed bounded convex set containing the conditional means, so all center functions belong to $L^2(\eta)$. The last expression is at least $q_M(\xi_s)$. Averaging the coins proves the assertion. No different continuation task is inserted into the loss, and no product structure of $\Lambda$ beyond the lower certificate is used.
\end{proof}

We next state a constructive upper certificate. It is convenient to allow a finite declared phase as in Section~\ref{sec:problem}; otherwise take one fiber. All constants below are uniform over phases and the specified strategies. A code has at most $M$ representatives in each phase. A cut lower must replace $M$ by the size of the full accessible tuple when its phase itself carries data.

Let $(S_v,d_v)$ be separable metric predictive fibers, with Borel envelope $w\ge1$. The update along a common report and physical action is $F_t$. A Borel relation $\mathcal A_v$ satisfies $x\in\mathcal A_v(x)$ and $w(z)\le w(x)$ for $z\in\mathcal A_v(x)$. It is fixed independently of the budget. Finite Borel covers have representatives $c_i^v$ and first-cell selectors $Q_M^v$ such that
\begin{equation}\label{eq:cover}
 Q_M^v(x)\in\mathcal A_v(x),\qquad
 d_v(x,Q_M^v(x))\le r_M w(x).
\end{equation}
For $s=km$, write $\mathscr F_{s,j}(z)$ for $j$ exact updates along the next physical word and the actions selected from that physical transcript. For every possible block-start past, true state $x$ and candidate $z$ in its named fiber, require one-step pathwise Lipschitz gain at most $L\ge1$ and
\begin{equation}\label{eq:pair}
 \E_x[d(\mathscr F_{s,m}(x),\mathscr F_{s,m}(z))^2\mid\mathcal H_s]
 \le\kappa d(x,z)^2,\qquad 0\le\kappa<1.
\end{equation}
The notation $d$ denotes the appropriate common fiber metric. For \emph{every} adapted sequence $z_0=z$, $v_j=F_{s+j}(z_{j-1},a_{s+j},Y_{s+j})$, $z_j\in\mathcal A(v_j)$, require
\begin{align}
 \E_x[w(z_m)^2\mid\mathcal H_s]&\le a w(z)^2+b,\label{eq:drift}\\
 \E_x[w(v_j)^2\mid\mathcal H_s]&\le A w(z)^2+B\quad(1\le j\le m),\label{eq:within}
\end{align}
where $0\le a<1$ and $b,A,B\ge0$. The conditional inequalities continue to hold when past independent algorithm coins are included. They quantify over arbitrary relation-preserving adapted perturbations, rather than assuming the performance of a preconstructed quantizer. Numerical implementations must still choose the same codebook and relation, with local error at most $r_Mw(v_j)+u$.

\begin{theorem}[Causal acquired-geometry and continuation transfer]\label{thm:transfer}
Fix a prepared experiment, a deterministic horizon and encoder-independent exploration. Suppose \eqref{eq:cover}--\eqref{eq:within} and a legal seed with $\E w(\widehat S_0)^2\le H_0$ and $\norm{d(S_0,\widehat S_0)}_2\le e_0$. Assume $P_T=f_T(S_T)$ with Lipschitz constant $L_f$ in the task norm. Put
\begin{align*}
 H&=\max\{H_0,b/(1-a)\},& H_*&=AH+B,\\
 \Lambda_j&=\sum_{i=0}^{j-1}L^i,& \Delta_M&=r_M\sqrt{H_*}+u,\\
 E_*&=\max\{e_0,\Lambda_m\Delta_M/(1-\sqrt\kappa)\},&
 E_M&=L^{m-1}E_*+\Lambda_{m-1}\Delta_M.
\end{align*}
For a one-fiber $M$-label model, or a model in which $M$ below counts the whole accessible tuple, define
\[
 W_M=\max\{q_M(\nu_T),\ \sup_{s,\,\text{certificates}}q_M(\xi_s)\}.
\]
All certificate measures are unnormalized: scaling $\underline\mu$ scales $q_M(\xi_s)$ by the same factor.
There is a causal predictor quantized after every raw call, with terminal readout error at most $v$ in root mean square, such that
\begin{equation}\label{eq:central}
 \begin{gathered}
 R_\cp(T,M)=B_T+q_M(\nu_T),\\
 B_T+W_M\le R_\on(T,M)
       \le B_T+(L_fE_M+v)^2.
 \end{gathered}
\end{equation}
For phase-indexed covers the upper uses the stated $M$ labels per phase and their charged full tuple; the lower always uses that tuple's cardinality.

The lower can also include every conditional continuation width of Theorem~\ref{thm:cutduality}. On the sequential-chart class of Theorem~\ref{thm:chart}, their supremum and the optimal recursively achievable excess are comparable; outside that class no sufficiency is claimed.

For any actual measure $\xi$, either $\nu_T$ or a dominated continuation measure, a submeasure of mass $\zeta>0$ with
\begin{equation}\label{eq:smallball}
 \sup_c\underline\xi(B(c,s_M))\le\zeta/(2M)
\end{equation}
contributes at least $\zeta s_M^2/2$ to the lower. Thus a matching actual terminal witness with $s_M\asymp r_M$, mass bounded below, $e_0,u,v=O(r_M)$ and fixed certificate constants gives matching checkpoint and online excess of order $r_M^2$. A larger continuation width instead certifies a separation from the terminal checkpoint curve. All terms refer to the one fixed task and actual preparation.
\end{theorem}
\begin{proof}
The exact selector in \eqref{eq:cover} is Borel and yields the recursion $\widehat S_t=Q_M(F_t(\widehat S_{t-1},a_t,Y_t))$. The numerical routine obeys the same relation. Equation~\eqref{eq:drift} proves inductively that the envelope second moment is at most $H$ at block endpoints; \eqref{eq:within} bounds each pre-rounding moment by $H_*$. Hence each local insertion of error has $L^2$ norm at most $\Delta_M$.

Start an exact shadow at $\widehat S_s$ for a block and feed it the same physical word. Telescoping the local insertions, with all subsequent pathwise gains included, gives
\[
 d(\mathscr F_{s,m}(\widehat S_s),\widehat S_{s+m})
 \le\sum_{j=1}^m L^{m-j}d(v_j,z_j).
\]
Conditional use of \eqref{eq:pair}, then Minkowski, shows that the root errors satisfy
\[
 e_{s+m}\le\sqrt\kappa e_s+\Lambda_m\Delta_M,
 \qquad e_{s+j}\le L^je_s+\Lambda_j\Delta_M\quad(0\le j<m).
\]
The first preserves $e_{km}\le E_*$ and the second gives $e_T\le E_M$. The shadow is a proof device: no block of reports, real posterior or second state is retained by the machine. Lipschitz readout, root-error addition and conditional projection prove the upper.

Lemma~\ref{lem:cut} gives all the lower bounds simultaneously for this same experiment. Their supremum remains a lower bound. Finally $M$ open balls satisfying \eqref{eq:smallball} cover at most half the submeasure. Outside their union the squared distance from every list is at least $s_M^2$, proving the stated contribution. This is valid in the continuation Hilbert space as well as in a finite-dimensional score space. The matching assertions follow directly; no ambient-dimensional volume has been assumed.
\end{proof}

\begin{corollary}[A necessary compatibility condition]\label{cor:necessary}
For a sequence of fixed-task experiments and budgets, an asserted bound
$R_\on(T,CM)-B_T\le C' q_M(\nu_T)$ implies
$W_{CM}\le C' q_M(\nu_T)$, with the same specified full-state cut. A family for which this inequality fails cannot have that claimed online/checkpoint matching, even if its terminal marginal has small covering number. Here $CM$ already denotes the entire data-dependent tuple at that cut. A separately stored phase of $Q$ values must be accounted for as $CMQ$ unless its value is deterministic at that cut; the constant $C$ cannot silently absorb a growing phase.
\end{corollary}
\begin{proof}
Apply the lower side of \eqref{eq:central}, or Lemma~\ref{lem:cut} alone when no stability certificate is available.
\end{proof}

\begin{example}[A quantitative delayed-query separation]\label{ex:query}
Prepare $d$ independent fair bits $X_1,\ldots,X_d$ and report them in $d$ paid calls. After the last retained cut report an independent uniform index $J\in\{1,\ldots,d\}$. The task is to predict $Y=X_J$ before its terminal audit. A checkpoint reading the full history after $J$ needs two labels for zero risk. If an online predictor had at most $M$ labels immediately before $J$, then
\begin{equation}\label{eq:randomaccess}
 R_\on\ge\tfrac12 h_2^{-1}\bigl((1-\log_2 M/d)_+\bigr),
\end{equation}
where $h_2^{-1}$ is the inverse binary entropy on $[0,1/2]$. In particular retaining at most a fixed fraction of $d$ bits gives a positive error bounded away from zero. Exact prediction requires $M\ge2^d$ and storing the bit string attains it.
\end{example}
\begin{proof}
The continuation measure is the law of $j\mapsto X_j$ in the uniform $L^2$ norm. For an independent shared seed $R$ and retained label $C$, $I(X;C\mid R)\le\log_2 M$. Independence of the bits and conditional subadditivity give
$\sum_j H(X_j\mid C,R)\ge d-\log_2M$.
Let $e$ be the average posterior bit-classification error. The two inequalities used are
\[
 h_2(e)\ge(1-\log_2M/d)_+,\qquad
 p(1-p)\ge p/2\quad(0\le p\le1/2).
\]
The first is concavity of binary entropy; the second bounds Bernoulli posterior variance by its classification error. Conditional projection proves \eqref{eq:randomaccess}. Zero error forces distinct positive-probability labels for all strings. This is the classical entropy mechanism behind random-access bounds~\cite{nayak}, used here to distinguish two information cuts of a single prediction task, not claimed as a new coding inequality.
\end{proof}

\begin{proposition}[Flagged violations of an admissibility relation]\label{prop:repair}
Suppose a finite implementation agrees with a routine satisfying the upper hypotheses until a flagged first failure, and, conditional on the current past and on no earlier flagged failure, its probability of failing at step $t$ is at most $\beta_t$. A repair routine using the same codebook can always be selected on a failure to satisfy the original relation and local-error bound. For a loss in $[0,L_*]$, the actual implementation has risk at most the repaired upper plus $L_*\min(1,\sum_{t\le T}\beta_t)$.
\end{proposition}
\begin{proof}
Drive the repaired routine by the same physical reports and coins; at a failure choose the certified selector instead. Encoder-independent exploration makes the physical law unchanged. The two retained trajectories agree until the first failure. A first-failure union bound gives its probability at most $\sum_t\beta_t$, and bounded loss gives the assertion. The reference estimate is unconditional; it is not incorrectly conditioned on the absence of failure. The proposition weakens implementation accuracy to a probabilistic certificate, but does not assert the existence of global relation-preserving covers from local data alone.
\end{proof}

\section{Two-sided continuation geometry}\label{sec:dual}
The preceding lower does not require the suffix to be independent of the prefix. Its strength nevertheless depends on finding a common sublaw. We now separate three questions: the exact one-cut compression problem, the production of nonzero sublaws from raw kernels, and recursive realization of their geometry. This leads to a quantitative converse and construction on a class with correlated reports. No asymptotic coding limit is taken.

\subsection{The exact cut problem and common channels}
Keep the fixed physical exploration and bounded task of Section~\ref{sec:problem}. Write $\mu=\law(H)$ and let $K(du\mid h)$ be the actual conditional law of the complete suffix. It includes every public action, phase, cost and report available before the decision. Public exploration randomness is part of the visible history. Private physical/controller randomness is integrated into $K$ unless actually reported; it is never made free decoder information.

Let $\mathcal D_M(\Lambda,z)$ be the least excess over $B_T$ when computation before and after one cut is unrestricted, but the only prefix-dependent object crossing it is a label in $[M]$. The decoder then receives the whole physical suffix. This is a relaxation of online prediction, not an assertion that a one-cut code can be run at every cut. Let $D\subset\R^q$ be a fixed compact convex set containing all task means. Decoder functions below are Borel and $D$-valued.

\begin{theorem}[One-cut comparison and conditional continuation]\label{thm:cutduality}
For any prepared standard-Borel experiment with encoder-independent exploration,
\begin{equation}\label{eq:exactfunctional}
 \mathcal D_M(\Lambda,z)=\inf_{a_1,\ldots,a_M}
 \int\min_{i\le M}\int\norm{z(h,u)-a_i(u)}^2K(du\mid h)\,\mu(dh).
\end{equation}
If finite measures $\mu_-,\mu_+$ and a suffix probability $\eta$ satisfy
\[
 \mu_-\otimes\eta\le\Lambda\le\mu_+\otimes\eta,
\]
then, for $Z(h)=z(h,\cdot)$ in $L^2(\eta;\R^q)$,
\begin{equation}\label{eq:cutsandwich}
 q_M(Z_*\mu_-)\le\mathcal D_M(\Lambda,z)
 \le q_M(Z_*\mu_+).
\end{equation}
In particular, $c\mu\otimes\eta\le\Lambda\le C\mu\otimes\eta$ gives comparison with $q_M(Z_*\mu)$ within factors $c,C$, for every $M$.

There is also a conditional lower when global product domination vanishes. Let $W$ be a standard-Borel physical variable, measurable from $H$, which is genuinely supplied again in $U$, or is explicitly given to the decoder as a relaxation. Write $\rho=\law(W)$ and disintegrate the actual joint law as $\Lambda_w$. Suppose jointly measurable kernels $\mu_{-,w}$ and $\eta_w$ satisfy $\mu_{-,w}\otimes\eta_w\le\Lambda_w$ for $\rho$-almost every $w$. Then
\begin{equation}\label{eq:conditionalcut}
 R_\on(T,M)-B_T\ge
 \int q_M\bigl((h\mapsto z(h,\cdot))_*\mu_{-,w};L^2(\eta_w)\bigr)\,\rho(dw).
\end{equation}
No normalization of the finite measures is intended. In particular $q_M(t\nu)=tq_M(\nu)$ for $t\ge0$. Conditioning in \eqref{eq:conditionalcut} is not permission to leave a continuous register uncharged in an upper construction.
\end{theorem}
\begin{proof}
Projection subtracts $B_T$. Fix the independent algorithm seed, including all future private decoder coins. Once its cut label is fixed, the entire subsequent finite sequence of Borel updates is a single Borel function of the complete suffix. Thus there are at most $M$ decoder functions. For a fixed list, integration against $K$ gives $M$ Borel distortions as functions of $h$; selecting the first minimum is Borel and optimal. Infimizing over lists proves \eqref{eq:exactfunctional}. Averaging independent seeds cannot improve its deterministic infimum. Projecting into $D$ cannot increase loss.

The lower comparison follows by integrating against the dominated product and then minimizing in $L^2(\eta)$. For the upper choose an approximating Hilbert center list for $Z_*\mu_+$, project its values into $D$, choose Borel representatives, and assign each prefix to its nearest center in that Hilbert norm. The actual distortion is bounded above by the dominating product distortion. Let the approximation error tend to zero. Neither comparison optimizes a policy-dependent law.

For the conditional statement, a code remains an $M$-function code after conditioning on $W=w$ and independent seeds. Apply the lower argument in each fiber and integrate. To verify measurability without assuming a measurable family of optimizers, fix a countable algebra generating the suffix Borel sigma field and a countable dense subset of $D$. The bounded simple functions on this algebra with those values are countable and dense in $L^2(\eta_w;D)$ for every $w$. Hence each fiber quantization value equals an infimum over countably many $M$-tuples of such functions. Each associated integral is a Borel function of $w$ by the kernel integration theorem. Their countable infimum is measurable. Conditional versions are immaterial because every dominated product is absolutely continuous with respect to its actual fiber law.
\end{proof}

Formula~\eqref{eq:exactfunctional} is the one-shot functional source-coding problem with decoder side information, written for the actual task conditional mean. We do not claim that ordinary encoder/decoder reduction as new. The useful implications are its comparison with acquired Hilbert measures, its conditional localization, and the recursive realization theorem below. Side information is genuinely causal here because it arrives after the retained-state cut; its coding problem after fixing the cut is an ordinary one-way problem.

\begin{lemma}[Raw constructions of common continuation laws]\label{lem:minorization}
The following give certificates for Theorem~\ref{thm:cutduality}.

\textup{(i)} If $K(du\mid h)$ has density $k(h,u)$ relative to a probability $\eta$ and $k(h,u)\ge\alpha(h)\ge0$ for $\mu\otimes\eta$-almost every $(h,u)$, then $\mu_-=\alpha\mu$. For a finite suffix alphabet with $\eta(u)>0$, one may take $\alpha(h)=\min_u K(u\mid h)/\eta(u)$. Upper density bounds give the upper comparison in the same way.

\textup{(ii)} Suppose a raw post-cut block, including \emph{all} its visible reports, costs and phases $R$, and its post-block hidden state $X'$, has conditional kernel
\begin{equation}\label{eq:regeneration}
 L(d r,dx'\mid h)\ge\alpha(h)\kappa(d r)\pi(dx').
\end{equation}
Suppose further that the subsequent controlled continuation has a kernel $F(du'\mid r,x')$ independent of the earlier prefix. This requires resetting or retaining in $x'$ every controller variable used later. Then the complete suffix $U=(R,U')$ has common sublaw
\[
 K(du\mid h)\ge\alpha(h)\eta(du),\qquad
 \eta(dr,du')=\kappa(dr)\int F(du'\mid r,x')\pi(dx').
\]
The same constructions hold after disintegration over $W$.
\end{lemma}
\begin{proof}
Multiply the density inequality by $\mu(dh)$ for (i). For (ii), integrate \eqref{eq:regeneration} against the nonnegative continuation kernel $F$; its normalization makes $\eta$ a probability. A hidden-state reset alone is insufficient if a prefix-dependent controller or an earlier visible block report is omitted. The displayed joint block minorization is precisely what prevents that omission. Integrating each fiber proves the final assertion.
\end{proof}

In (ii) a Doeblin block of a prepared instrument produces the product certificate before any geometric estimate is made. Its mass is $\int\alpha\,d\mu$, and this mass multiplies the squared-distortion lower. Conditioning on a rare successful block without its mass would be a different experiment.

\begin{example}[A diagonal interface with no global product certificate]\label{ex:diagonal}
Prepare an independent uniform $W\in[0,1]$ and $d$ fair bits $X_1,\ldots,X_d$. Report $W$ and then the bits by paid calls. After the cut report $(W,J)$, where $J$ is independent uniform on $[d]$, and score $Y=X_J$ after the decision. Every call, including the reissued $W$, is physical and charged. The actual prefix/suffix law is supported on equality of its two $W$ coordinates and admits no nonzero product submeasure. Indeed a nonzero independent product supported on the diagonal must have both $W$ marginals concentrated at one common point, whereas every prefix submeasure has a nonatomic $W$ marginal.

Nevertheless, conditional on $W=w$ the continuation response is the bit vector in $L^2(\mathrm{Unif}[d])$. Theorem~\ref{thm:cutduality} gives exactly the random-access obstruction of Example~\ref{ex:query}; exact prediction still requires $M\ge2^d$, while the final checkpoint needs only two labels. Thus failure of a global certificate neither shows easy compression nor invalidates conditional acquired geometry. No copy of $W$ is retained for free by the algorithm: it is reissued by the specified physical experiment.
\end{example}

\subsection{From continuation charts to a finite recursion}
We give an explicit class where the earlier continuation geometry is both necessary and sufficient in order. Its acquisition is sequential but need not be independent or stationary. Fix $d\ge1$. A prepared experiment first reports $H_1,\ldots,H_d\in\R$ in $d$ paid calls, in that order, with any jointly specified law $\mu$. It next reports a single finite-dimensional packet $U$ by a paid call, and scores a scalar $Y\in[-1,1]$ after the decision. Preparation and the scored audit are also charged. There are exactly $d+3$ raw calls. Report and phase registers are not rereadable across cuts. The only admitted control in this theorem is this acquisition order; the theorem is not an optimization of that order.

Assume that $z(h,u)=\E[Y\mid H=h,U=u]$ has a bounded Borel extension to all $h\in\R^d$ and satisfies
\begin{equation}\label{eq:chartlip}
 |z(h,u)-z(h',u)|\le L\norm{h-h'}_2
 \quad\text{for every }h,h',u.
\end{equation}
Assume $\E(1+|H_j|)^{2+\zeta}\le A$ for some $\zeta>0$, uniformly in $j$. Finally there is a compact Borel set $E\subset\R^d$ with $\mu(E)=p>0$, $\mu|_E$ having Lebesgue density at most $b<\infty$, and a suffix probability $\eta$ such that
\begin{equation}\label{eq:chartlower}
 K(du\mid h)\ge\alpha\eta(du)\ (h\in E),\qquad
 \norm{z(h,\cdot)-z(h',\cdot)}_{L^2(\eta)}
 \ge l\norm{h-h'}_2\ (h,h'\in E),
\end{equation}
where $\alpha,l>0$. These are hypotheses on the law and task response computed from the raw kernels. No covariance, Fisher metric or pre-existing manifold is supplied as geometric input. The second inequality identifies locally distinguishable acquired directions through actual future tests.

\begin{theorem}[Continuation-chart realization]\label{thm:chart}
Under \eqref{eq:chartlip}--\eqref{eq:chartlower}, put
$\xi=(h\mapsto z(h,\cdot))_*(\alpha\mu|_E)$ and let $B$ be the full-history Bayes risk for this one task. There are constants $0<c\le C<\infty$, depending only on $d,\zeta,A,L,p,b,\alpha,l$, such that for every $M\ge1$,
\begin{equation}\label{eq:chartlaw}
 cM^{-2/d}\le q_M(\xi)\le R_\on(M)-B\le CM^{-2/d}.
\end{equation}
The upper uses a genuine $M$-label recursion after each scalar report and after the suffix packet. Its constants do not require independence of the prefix coordinates. If $W_M^{\rm all}$ denotes the maximum of terminal quantization and all the unconditional and conditional cut lower bounds above, then
\begin{equation}\label{eq:chartcomplete}
 R_\on(M)-B\asymp W_M^{\rm all}\asymp q_M(\xi)\asymp M^{-2/d}.
\end{equation}
Thus these cut widths are quantitatively sufficient for this declared sequential-chart class, not for all causal experiments. The phase has $d+2$ pre-audit values; if stored, it multiplies the persistent count. The scalar task always has $R_\cp(M)-B\le M^{-2}$.
\end{theorem}
\begin{proof}
A Hilbert ball of radius $r$ meeting the image of $E$ has preimage contained in a Euclidean ball of radius $2r/l$, by \eqref{eq:chartlower}. Its $\xi$-mass is at most $\alpha b v_d(2r/l)^d$, where $v_d$ is the volume of the Euclidean unit ball. Taking
\[
 r=\frac l2\left(\frac{p}{2b v_d M}\right)^{1/d}
\]
shows that $M$ such balls cover at most half the actual mass $\alpha p$. The remaining mass contributes at least $\alpha p r^2/2$. Theorem~\ref{thm:cutduality} proves the first two inequalities of \eqref{eq:chartlaw}.

For the construction we give the finite scalar quantizer explicitly. Put $a_0=\zeta/2$ and, for integer $k\ge1$, map $x\ge0$ to
\[
 t(x)=1-(1+x)^{-a_0},\quad j=\lfloor kt(x)\rfloor,\quad
 Q_k(x)=(1-j/k)^{-1/a_0}-1;
\]
extend oddly to negative $x$. It has at most $2k-1$ values. Since $0\le Q_k(x)\le x$ and $t'$ decreases,
\begin{equation}\label{eq:compandpoint}
 |x-Q_k(x)|\le\frac{(1+|x|)^{1+a_0}}{a_0k},\qquad
 \E|H_j-Q_k(H_j)|^2\le\frac{A}{a_0^2k^2}.
\end{equation}
This includes the two unbounded outside cells; truncating a Gaussian without paying its tails is unnecessary.

Let $m=\lfloor M^{1/d}\rfloor$ and $k=\lfloor(m+1)/2\rfloor$. On receiving $H_j$ append its quantizer index. After $j$ reports there are at most $m^j\le M$ possible tuples. No actual real prefix is retained. At the packet call compute $z(\widehat h,U)$ from this finite tuple and the current packet, then retain its nearest midpoint in an $M$-cell uniform partition of $[-1,1]$. The old tuple is discarded. The packet and all arithmetic workspace are erased before the next cut. Equation~\eqref{eq:chartlip}, the marginal estimates \eqref{eq:compandpoint}, and Minkowski give root excess at most
\[
 L\sqrt{dA}/(a_0k)+M^{-1}\le C_0M^{-1/d}.
\]
Projection onto the full-history conditional mean proves the upper. This is a finite-label Borel implementation, not yet a finite-bit implementation of arbitrary real coefficients; those resources are discussed below.

Each cut lower is at most the risk of every online predictor, while the particular $\xi$ lower is at least $cM^{-2/d}$. The proved upper therefore yields \eqref{eq:chartcomplete}, with no assumption that an optimal one-cut center list is recursively executable. At the final checkpoint uniform scalar quantization of $z(H,U)\in[-1,1]$ has error at most $M^{-2}$.
\end{proof}

The hypotheses can be checked without a global smooth parameterization. Only one actual positive-mass chart is needed for the lower, while the marginal moment and global response modulus control all acquired tails for the upper. The hard intermediate memory constraint is resolved by coordinatewise acquisition, not by an assumed contractive posterior. Degenerating chart mass or response rank changes the constants and is not hidden in an ambient dimension. The next theorem verifies every ingredient in a genuinely noisy, correlated continuation experiment.

\section{The acquired-dimension transfer theorem}\label{sec:serial}
A cut obstruction and a recursively selectable cover measure different objects unless they can be compared through the same experiment. The following result supplies that comparison along an entire finite experiment. It permits changes of dimension, noninteger geometric exponents, dependent reports and expansive updates. Its constants retain the length and gains of the experiment. A finite-horizon statement is not converted into a time-uniform stability theorem.

Fix the preparation, encoder-independent exploration and bounded terminal task of Section~\ref{sec:problem}, and a deterministic number $T\ge1$ of predecision reports. Write $O_t$ for the complete report at call $t$, including its visible action, phase and resource record. All full-history conditional means use the actual joint law. Preparation and audit are additional charged calls. Algorithm seeds are independent of this law.

A \emph{serial task realization} consists of standard-Borel metric spaces $(S_t,d_t)$, jointly Borel maps $F_t:S_{t-1}\times\mathcal O_t\to S_t$, and a known initial point $z_0$, such that
\begin{equation}\label{eq:serialstate}
 Z_0=z_0,\qquad Z_t=F_t(Z_{t-1},O_t),\qquad Z_T=P_T.
\end{equation}
Here $S_T$ is a bounded convex subset of the task's Euclidean space and $d_T$ is its task norm. A task realization may be a proper factor of the full predictive quotient. Its sufficiency for this scored task does not assert that it generates the entire future report law. The latter assertion is required separately when an experiment morphism uses it.

The maps and spaces in \eqref{eq:serialstate} are verified from the prepared kernels and future tests. They are not an assumed parameter manifold. We use the following quantitative hypotheses, fixed independently of $M$.

\paragraph{Global recursion and covers.}
There are finite constants $L_t,a_t>0$, $\lambda_t\ge0$, Borel functions $V_t:S_t\to[1,\infty)$ and $b_t:\mathcal O_t\to[0,\infty)$, with $\|b_t(O_t)\|_2<\infty$, such that for every pair of candidate states and every report in the declared complete domain,
\begin{align}
 d_t(F_t(x,o),F_t(x',o))&\le L_t d_{t-1}(x,x'),\label{eq:seriallip}\\
 V_t(F_t(x,o))&\le b_t(o)+\lambda_t V_{t-1}(x).\label{eq:seriallyap}
\end{align}
For each integer $M\ge1$ there is a Borel map $Q_{t,M}:S_t\to S_t$ with at most $M$ values satisfying
\begin{equation}\label{eq:serialcover}
 d_t(x,Q_{t,M}x)\le a_t M^{-1/s_t}V_t(x),\qquad
 V_t(Q_{t,M}x)\le V_t(x),
\end{equation}
where $s_t>0$. Singleton stages are allowed: their quantizer is the identity, their insertion radius is zero, and they are assigned $s_t=0$. They are omitted from expressions containing $1/s_t$. These are global candidate-domain conditions, not a cover only of the true trajectory. The Lyapunov inequality quantifies over candidates before choosing a code.

\paragraph{Actual continuation mass.}
At every nonsingleton cut $t$, let $H_t$ be the complete acquired prefix and let $U_t$ contain all remaining predecision visible information. A probability $\eta_t$ and a finite measure $\underline\mu_t$ on prefixes satisfy
\begin{equation}\label{eq:serialdomination}
 \underline\mu_t\otimes\eta_t\le\law(H_t,U_t).
\end{equation}
On this product a bounded Borel conditional-mean version factors as
\[
 \E[Y\mid H_t=h,U_t=u]=g_t(Z_t(h),u).
\]
Put $\rho_t=(Z_t)_*\underline\mu_t$ and $p_t=\rho_t(S_t)>0$. On a Borel set carrying $\rho_t$, assume
\begin{align}
 \rho_t(B_{d_t}(x,r))&\le C_t r^{s_t}\quad(x\in S_t,\ r>0),\label{eq:serialmass}\\
 \|g_t(x,\cdot)-g_t(x',\cdot)\|_{L^2(\eta_t)}
 &\ge \ell_t d_t(x,x')\qquad(\ell_t>0).\label{eq:serialobservable}
\end{align}
The response functions are bounded and jointly Borel on the declared domains, so their $L^2$ evaluation images are measurable in separable Hilbert spaces. At the terminal cut $U_T$ is a point and $g_T(x)=x$. The finite measure in \eqref{eq:serialdomination} retains its actual mass; a success event is not renormalized. The exponent $s_t$ is witnessed both by an actual observable submeasure and by a global selectable cover. Equality of a formal ambient dimension with $s_t$ is neither assumed nor used.

\begin{theorem}[Serial acquired-geometry and compatibility]\label{thm:serial}
Suppose \eqref{eq:serialstate}--\eqref{eq:serialobservable} hold, with $s_T>0$, and set $s_*:=\max_{1\le t\le T}s_t$. Let
\[
 A_0=V_0(z_0),\qquad A_t=\|b_t(O_t)\|_2+\lambda_t A_{t-1},\qquad
 \Gamma_t=a_tA_t\prod_{j=t+1}^{T}L_j.
\]
At each nonsingleton cut let
$\xi_t=(x\mapsto g_t(x,\cdot))_*\rho_t$ in $L^2(\eta_t)$.
Then, for every integer $M\ge1$,
\begin{equation}\label{eq:serialexplicit}
 \max_t q_M(\xi_t)\le R_\on(T,M)-B_T
 \le \left(\sum_{t:s_t>0}\Gamma_tM^{-1/s_t}\right)^2.
\end{equation}
There are $0<c\le C<\infty$, depending only on the displayed certificates and the finite experiment, such that
\begin{align}
 R_\on(T,M)-B_T&\asymp \max_t q_M(\xi_t)\asymp W_M^{\rm all}
                    \asymp M^{-2/s_*},\label{eq:serialrate}\\
 R_\cp(T,M)-B_T&=q_M(\nu_T)\asymp M^{-2/s_T}.\label{eq:serialcheckpoint}
\end{align}
Here $W_M^{\rm all}$ is the supremum of all valid same-task cut lower bounds, including terminal quantization and conditional cuts when available. Consequently checkpoint and online excesses match within constant factors in this class if and only if $s_*=s_T$. Otherwise their ratio is of order
\begin{equation}\label{eq:serialratio}
 M^{\,2/s_T-2/s_*}.
\end{equation}
The upper is one Borel predictor with at most $M$ data-dependent labels at every cut, not a sequence of unrelated optimal checkpoint codes. A stored deterministic phase has at most $T+1$ values and multiplies the full persistent alphabet. Any data-dependent controller or simulator state is additional charged state, not such a phase.
\end{theorem}
\begin{proof}
For the lower, a Hilbert ball of radius $r$ meeting the image of the carrying set of $\rho_t$ has preimage contained in a $d_t$-ball of radius $2r/\ell_t$: choose one preimage in the intersection and use \eqref{eq:serialobservable} and the triangle inequality. Hence
\[
 \xi_t(B(c,r))\le C_t(2r/\ell_t)^{s_t}.
\]
With $r=(\ell_t/2)(p_t/(2C_tM))^{1/s_t}$, $M$ such balls cover at most $p_t/2$. It follows that
\begin{equation}\label{eq:seriallowerconstant}
 q_M(\xi_t)\ge \frac{p_t\ell_t^2}{8}
                 \left(\frac{p_t}{2C_tM}\right)^{2/s_t}.
\end{equation}
The actual product domination and Lemma~\ref{lem:cut} transfer this lower to every online predictor. The argument includes independent randomized algorithms because fixing all their coins leaves at most $M$ Borel suffix-response functions. No information from unreported physical coins is supplied to the decoder.

For the upper, construct
\[
 \widehat Z_0=z_0,\quad \widetilde Z_t=F_t(\widehat Z_{t-1},O_t),\quad
 \widehat Z_t=Q_{t,M}(\widetilde Z_t).
\]
Only the index of $\widehat Z_t$ is retained. The current report, the previous index and arithmetic workspace are discarded at the cut. By \eqref{eq:seriallyap} and the inward inequality in \eqref{eq:serialcover}, induction and Minkowski give
\[
 \|V_t(\widetilde Z_t)\|_2\le A_t,\qquad
 \|V_t(\widehat Z_t)\|_2\le A_t.
\]
This bound uses no independence between reports and candidates. If
$e_t=\|d_t(Z_t,\widehat Z_t)\|_2$, then
\[
 e_t\le L_t e_{t-1}+a_tA_tM^{-1/s_t},\qquad e_0=0.
\]
At singleton stages the insertion is zero. Unrolling this deterministic inequality gives the upper in \eqref{eq:serialexplicit}. Conditional projection onto $P_T=Z_T$ makes $e_T^2$ precisely the excess risk of this predictor.

Choose a cut with $s_t=s_*$. Equation~\eqref{eq:seriallowerconstant} gives $cM^{-2/s_*}$. Since $M^{-1/s_t}\le M^{-1/s_*}$ at every active cut, the proved upper is at most $(\sum_t\Gamma_t)^2M^{-2/s_*}$. Every valid cut lower is no larger than every online risk, so taking its supremum proves all comparisons in \eqref{eq:serialrate}; no optimal codebook is assumed to be causally selectable.

For the checkpoint upper apply $Q_{T,M}$ directly to $P_T$. The exact recursion, without any quantizers, obeys the same Lyapunov bound, so its mean squared insertion is at most $a_T^2A_T^2M^{-2/s_T}$. The terminal case of \eqref{eq:seriallowerconstant} gives the reverse inequality for $q_M(\nu_T)$. The checkpoint projection identity proves \eqref{eq:serialcheckpoint}. Dividing the two matched bounds gives \eqref{eq:serialratio} and the stated criterion.
\end{proof}

The theorem is a quantitative sufficiency statement for the full family of continuation cuts on this certificate class. It does not say that cut lower bounds alone construct a recursion in every experiment. Its extra content is the raw-verifiable global update and covering mechanism, including a noncircular candidate-moment bound. Neither a single smooth chart nor contraction at every step is required. The block theorem of Section~\ref{sec:transfer} supplies stronger, time-uniform estimates when the raw experiment has its stated block stability.

\begin{lemma}[Unbounded Euclidean acquisition and changing rank]\label{lem:serialcompander}
The global hypotheses of Theorem~\ref{thm:serial} hold for a finite sequence of Euclidean candidate spaces $S_t=\R^{r_t}$, with $s_t=r_t$, whenever the raw update maps are globally Lipschitz in the preceding state and
\[
 \|F_t(x,o)\|_2\le L_t\|x\|_2+\beta_t(o),\qquad
 \E(1+\beta_t(O_t))^{2+\zeta}<\infty
\]
for some $\zeta>0$. Coordinates may be dependent and $r_t$ may increase or decrease. A bounded terminal task is handled by a separate bounded cover in its actual affine or metric domain. If its actual continuation maps have the positive-mass observable charts of \eqref{eq:serialdomination}--\eqref{eq:serialobservable}, then the matching conclusions hold with these exponents.
\end{lemma}
\begin{proof}
Put $q=1+\zeta/2$ and $V_t(x)=(1+\|x\|_2)^q$. The scalar inward compander constructed in \eqref{eq:compandpoint}, applied coordinatewise with
$m=\lfloor M^{1/r_t}\rfloor$ and $k=\lfloor(m+1)/2\rfloor$, has at most $(2k-1)^{r_t}\le M$ representatives. It does not increase Euclidean norm, and its pointwise error is at most
\[
 \frac{4\sqrt{r_t}}{\zeta/2}M^{-1/r_t}(1+\|x\|_2)^q.
\]
Here $k\ge M^{1/r_t}/4$, including $M=1$. Thus the two cover inequalities hold globally, including the unbounded outside cells. The elementary inequality $(a+b)^q\le2^{q-1}(a^q+b^q)$ gives \eqref{eq:seriallyap} with
\[
 b_t(o)=2^{q-1}(1+\beta_t(o))^q,\qquad
 \lambda_t=2^{q-1}L_t^q.
\]
The moment assumption makes $b_t(O_t)$ square integrable. This proves the required candidate bound before constructing the eventual code. Bounded compact metric domains with a global $CM^{-1/s_t}$ cover instead use $V_t=1$, $b_t=1$, $\lambda_t=0$. First-nearest selection from a finite net is Borel. These observations also give the requisite bounded terminal covers.
\end{proof}

\begin{corollary}[One-metric numerical and causal composition]\label{cor:serialerrors}
Replace $F_t$ in the implemented recursion by a Borel map $\widetilde F_t$ into the same candidate domain. Suppose it obeys the same Lyapunov envelope (or a displayed enlarged one) and, for every admitted adapted candidate process, its root error relative to $F_t$ under the actual experiment is at most $u_t$. Then
\begin{equation}\label{eq:serialerrorprofile}
 R_{\rm implemented}-B_T\le
 \left[\sum_t\left(a_tA_tM^{-1/s_t}+u_t\right)
                    \prod_{j=t+1}^TL_j\right]^2,
\end{equation}
with zero cover insertion at singleton stages. A calibration displacement bounded by $\delta_t$ contributes at most $\kappa_t\delta_t$ to $u_t$ when the raw update has a report-modulus certificate $\kappa_t$ in the same state metric. Numerical evaluation errors add to this root error, not to an unrelated task metric.

Under a parameter-independent causal morphism of Theorem~\ref{thm:morphism}, with complete-law defect $\varepsilon$ and loss bounded by $L_{\max}$, the transferred \emph{absolute} risk upper gains at most $L_{\max}\varepsilon$. Predictor, controller, simulator and stored phase alphabets multiply. Read-only program length, scratch, numerical output precision, physical calls and time remain separate resources. No allowed error budget is asserted to be an attained lower-risk floor.
\end{corollary}
\begin{proof}
The same Lyapunov induction bounds the perturbed prequantization states. Triangle inequality adds $u_t$ to the recurrence for $e_t$; unroll it and use the terminal projection identity. A coupled perturbed report contributes $\kappa_t\delta_t$ by its stipulated report modulus. The complete-law total-variation inequality for a loss in $[0,L_{\max}]$ gives the last assertion after the task error has been estimated. It compares one scored target law; it does not subtract two different Bayes baselines. The serial resource tuple is exactly that of the typed morphism theorem.
\end{proof}

\section{Noisy continuation and a sharp causal separation}\label{sec:noisy}
We now compute the earlier continuation width in a noisy regression experiment. The future report contains both a noisy linear observation and a continuously varying query. Responses overlap, and no bit or query index can be recovered exactly. The same bounded binary task is used for the checkpoint and the online predictor.

Fix $d\ge2$, $\tau,\sigma>0$, and a nonzero known $a\in\R^d$. Prepare $\Theta\sim N(0,I_d)$. In $d$ successive paid calls report
\begin{equation}\label{eq:noisyprefix}
 H_j=\Theta_j+\tau Z_j,\qquad 1\le j\le d,
\end{equation}
where $Z_1,\ldots,Z_d$ are independent standard normals. The next paid call reports the single packet
\begin{equation}\label{eq:noisypacket}
 U=(V,B),\qquad V=a^\top\Theta+\sigma Z_0,
 \qquad B\sim\mathrm{Unif}(\mathbb S^{d-1}),
\end{equation}
with $Z_0,B$ independent of the preceding preparation. After the decision a paid audit reports a binary $Y$ with conditional success probability
\[
 \psi(B^\top\Theta),\qquad \psi(t)=\tfrac12(1+\tanh t).
\]
This audit is not available for prediction. Preparation, reports and audit cost $N=d+3$ raw calls. The legal strategy class consists of this fixed acquisition order and all randomized Borel prediction recursions subject to the declared retained-state bound. In particular the packet cannot be split into free later calls, and the prediction rule cannot change the physical exploration. Squared loss is $|Y-\widehat p|^2$, with $\widehat p\in[0,1]$.

\begin{theorem}[Noisy regression continuation]\label{thm:noisy}
For the raw experiment \eqref{eq:noisyprefix}--\eqref{eq:noisypacket}, let
\[
 P=\E[Y\mid H,V,B],\quad B_*=\E[P(1-P)],\quad \nu=\law(P).
\]
There is an actual dominated continuation measure $\xi$ at the cut after $H_d$ and constants $0<c<C<\infty$, depending only on $d,\tau,\sigma,a$, such that for all $M\ge1$,
\begin{align}
 cM^{-2/d}&\le q_M(\xi)\le R_\on(M)-B_*\le CM^{-2/d},\label{eq:noisylaw}\\
 cM^{-2}&\le R_\cp(M)-B_*=q_M(\nu)\le CM^{-2}.\label{eq:noisycp}
\end{align}
Thus earlier continuation geometry, rather than the scalar terminal acquired law, is essential to the online resource exponent. All intermediate data-dependent predictor state is included in $M$. A stored deterministic phase is separately charged. Every normal noise variance is strictly positive, the query is continuous, and the conditional suffix law is genuinely prefix dependent.

The same conclusion holds for any query probability $\lambda$ on the unit sphere satisfying $\int bb^\top\lambda(db)\succeq gI_d$ with $g>0$, with constants also depending on $g$. One may also replace the prefix by $H=G\Theta+\tau Z$ for any fixed invertible matrix $G$, reporting its coordinates in the same paid order; constants may then depend on $G$. No whitening operation on an unretained exact prefix is allowed. The response-rank condition is therefore a condition on executable future tests, not on a formal parameter dimension.
\end{theorem}
\begin{proof}
We verify the conditions of Theorem~\ref{thm:chart} from the normal densities. Multiplying the prior density by those of \eqref{eq:noisyprefix} and completing squares gives
\[
 \Theta\mid H=h\sim N(ch,v_0I_d),\qquad
 c=(1+\tau^2)^{-1},\quad v_0=\tau^2/(1+\tau^2).
\]
Multiplication by the likelihood of $V=v$ and a second completion of squares give
\begin{equation}\label{eq:posteriornoisy}
 \Theta\mid H=h,V=v\sim N(Dch+kv,\Sigma),
 \quad k=\frac{v_0a}{s_V^2},\quad D=I_d-ka^\top,
 \quad \Sigma=v_0I_d-\frac{v_0^2aa^\top}{s_V^2},
\end{equation}
where $s_V^2=\sigma^2+v_0\norm a_2^2$. Both $D$ and $\Sigma$ are positive definite; $\lambda_{\min}(D)=\sigma^2/s_V^2$. Thus these matrices are derived conditional quantities, not geometric assumptions.

For $s\ge0$ put $g_s(t)=\E\psi(t+sZ)$, $Z\sim N(0,1)$. The actual response function is
\begin{equation}\label{eq:regresponse}
 z(h,v,b)=g_{s_b}\bigl(b^\top(Dch+kv)\bigr),\qquad
 s_b^2=b^\top\Sigma b.
\end{equation}
It is Borel and lies in $[0,1]$. Dominated differentiation gives
\begin{equation}\label{eq:smoothlower}
 0<g_s'(t)=\tfrac12\E\operatorname{sech}^2(t+sZ)\le\tfrac12.
\end{equation}
Consequently the global prefix Lipschitz constant is $c\norm D_{\mathrm{op}}/2$, independent of the unbounded value of $v$. The acquired $H$ has density $N(0,(1+\tau^2)I_d)$, so every marginal moment needed in Theorem~\ref{thm:chart} is finite.

Fix any $R,T>0$ and use the actual cube $E=[-R,R]^d$. Its acquired mass $p=\PP(H\in E)$ is positive and its unnormalized density is bounded above. From the same density calculation,
\[
 V\mid H=h\sim N(a^\top ch,s_V^2).
\]
For $h\in E$ and $|v|\le T$ this density has a strictly positive minimum $m_0$. With
\[
 \eta(dv,db)=\mathrm{Unif}[-T,T](dv)\lambda(db),\qquad
 \alpha=2Tm_0>0,
\]
we obtain $K(du\mid h)\ge\alpha\eta(du)$ on $E$. This is the raw common-density construction of Lemma~\ref{lem:minorization}; neither $p$ nor $\alpha$ is normalized away. It also quantifies report continuity:
\begin{equation}\label{eq:noisytv}
 \norm{K(\cdot\mid h)-K(\cdot\mid h')}_{\TV}
 \le \frac{c\norm a_2}{\sqrt{2\pi}\,s_V}\norm{h-h'}_2.
\end{equation}
Indeed equal-variance scalar normal total variation is at most the mean displacement divided by $\sqrt{2\pi}s_V$; this follows by integrating the absolute derivative of the translated normal density. The independent query kernel does not increase it.

For $h\in E$, $|v|\le T$ and unit $b$, the argument of $g_{s_b}$ has absolute value at most
\[
 L_0=c\norm D_{\mathrm{op}}\sqrt d R+\norm k_2T,
 \qquad 0<s_b\le\sqrt{v_0}.
\]
On this interval \eqref{eq:smoothlower} has the explicit uniform lower
\[
 \gamma=\tfrac12\PP(|Z|\le1)
               \operatorname{sech}^2(L_0+\sqrt{v_0})>0.
\]
The mean value theorem, using the segment joining $h,h'\in E$, therefore gives
\begin{align*}
 \norm{z(h,\cdot)-z(h',\cdot)}_{L^2(\eta)}^2
 &\ge\gamma^2 c^2\int |b^\top D(h-h')|^2\lambda(db)\\
 &\ge\gamma^2c^2g\lambda_{\min}(D)^2\norm{h-h'}_2^2.
\end{align*}
For the uniform sphere, $g=1/d$. Thus the response functions form a co-Lipschitz acquired chart even though each fixed response is scalar and noisy. Put $\xi=(h\mapsto z(h,\cdot))_*(\alpha\mu|_E)$. Theorem~\ref{thm:chart} now proves \eqref{eq:noisylaw}, including a recursion that quantizes each actual report as it arrives. It does not store the exact posterior vector.

It remains to identify the terminal checkpoint exponent in the same task. The posterior mean $m'=D c H+kV$ is centered normal with covariance $C'=I_d-\Sigma$. To see this without a matrix identity shortcut, use total covariance in \eqref{eq:posteriornoisy}: the conditional covariance is the constant $\Sigma$, while the prior covariance is $I_d$. Since $0\prec\Sigma\preceq v_0I_d\prec I_d$, the eigenvalues of $C'$ lie in a fixed compact subset of $(0,1)$. Conditional on $B=b$, the scalar $T_b=b^\top m'$ is normal with variance $b^\top C'b$ bounded uniformly away from zero and infinity. Fix $A_0>0$ and retain the actual subevent $|T_b|\le A_0$. Its mass is uniformly positive. On it the map $t\mapsto g_{s_b}(t)$ is strictly increasing with derivative bounded below by a positive constant independent of $b$, by the preceding derivative estimate. The change-of-variable formula thus bounds the density of this scalar posterior submeasure uniformly above. Mixing over $b$ preserves both a positive submeasure mass and this density bound. It follows by the one-dimensional small-ball argument that $q_M(\nu)\ge cM^{-2}$. Uniform midpoint quantization on $[0,1]$ gives the reverse bound $q_M(\nu)\le1/(4M^2)$. The projection identity uses exactly $B_*$ on both sides and proves \eqref{eq:noisycp}.

For the stated correlated-prefix extension, direct square completion gives $m=C_Hh$ and $\Sigma_0=(I_d+\tau^{-2}G^\top G)^{-1}$, where $C_H=\Sigma_0\tau^{-2}G^\top$ is invertible. Put $s_V^2=\sigma^2+a^\top\Sigma_0a$, $k=\Sigma_0a/s_V^2$, $D=I_d-ka^\top$, and $\Sigma=\Sigma_0-\Sigma_0aa^\top\Sigma_0/s_V^2$. Now $D$ is invertible, with determinant $\sigma^2/s_V^2$, and $0\prec\Sigma\prec I_d$. Replace $cD$ by $DC_H$ in the global modulus and use its positive least singular value in the response lower. The acquired prefix is a nondegenerate, generally correlated Gaussian, so the same compact mass, density and marginal moment estimates apply. In the report-continuity estimate replace $c\norm a_2$ by $\norm{C_H^\top a}_2$. The checkpoint proof uses $I_d-\Sigma$ unchanged. Coordinate companding still acts on the actually reported $H_j$, not on inaccessible exact whitened values. This verifies every step for the extension.
\end{proof}

\begin{corollary}[Stable numerical and simulation composition]\label{cor:noisynumeric}
In Theorem~\ref{thm:noisy}, suppose the final response computation, using the retained representative and the current packet, differs from its exact value by an $L^2$ amount at most $u$. Suppose each acquired prefix report is causally perturbed by at most $\delta$ before quantization. A label-preserving implementation then has
\begin{equation}\label{eq:noisynumeric}
 R_\on-B_*\le C\bigl(M^{-1/d}+\delta+u\bigr)^2.
\end{equation}
If this complete implemented experiment and bounded scored task are transferred through a causal simulator with complete-law defect at most $\varepsilon$, the absolute risk upper gains at most $\varepsilon$. The simulator, its phase and its controller states must be multiplied into the full retained alphabet as in Theorem~\ref{thm:morphism}. Neither an allowed error $u$ nor an allowed defect $\varepsilon$ is asserted to be an unavoidable risk floor.
\end{corollary}
\begin{proof}
For a perturbed report $\widetilde H_j$ with $|\widetilde H_j-H_j|\le\delta\le1$, its $(2+\zeta)$ moment is bounded by a fixed multiple of that of $H_j$. Apply the explicit compander of \eqref{eq:compandpoint} to $\widetilde H_j$. Minkowski gives $\norm{H-\widehat H}_{L^2}\le C_1M^{-1/d}+\sqrt d\delta$. The global response Lipschitz estimate, final midpoint quantization and the evaluation error yield \eqref{eq:noisynumeric}. For $\delta>1$ the bounded loss gives the same conclusion after increasing $C$. Clipping computed probabilities into $[0,1]$ is nonexpansive. The total variation inequality for a loss in $[0,1]$ proves the final assertion for the absolute risk, without incorrectly transferring a Bayes baseline for free.
\end{proof}

\paragraph{Computation and scope.}
The finite-state theorem uses exact Borel report comparisons and readout; it does not grant a free infinite-precision implementation. For example, taking $\zeta=2$ makes every finite compander threshold and representative rational. With a specified finite-precision report interface, normal integral evaluations with a declared error $u$, and a certified coefficient table, the corollary supplies a quantitative error target. The storage of that table, scratch, output precision and time are additional resources, not contained in $\log_2 M$. The complete finite-bit feasible accounting theorem below is proved for the progressive experiment. We do not claim a matching bit-operation lower for this regression model.

The example is a fixed-acquisition Bayesian experiment, not an unknown-kernel learning theorem or an optimal-control theorem. Its purpose is different: a positive-noise continuation channel and a continuum of overlapping nonlinear future tests force a strict online/checkpoint exponent separation, and the general continuation-chart theorem is the common lower and construction mechanism. A final scalar posterior density alone cannot give the online exponent in \eqref{eq:noisylaw}.

\section{Observable rank and singular acquisition}\label{sec:rank-singu}
The exponent in Theorem~\ref{thm:serial} is determined by acquired response directions at each cut. The following two verifications distinguish it both from the dimension of an ambient hidden variable and from the dimension of the final forecast. The first uses noisy correlated Gaussian reports; the second has a genuinely singular acquired distribution and a noisy suffix. Neither is used to prove the general theorem.

\subsection{Noisy regression with a changing observable rank}
Retain the raw regression experiment of Section~\ref{sec:noisy}, allowing the invertible matrix $G$ in the prefix $H=G\Theta+\tau Z$. Change only the preparation of the query: its probability $\lambda$ is supported on the unit sphere of a linear subspace $L\subset\R^d$ of dimension $1\le r\le d$. Assume
\begin{equation}\label{eq:rankframe}
 \int bb^\top\lambda(db)\succeq gP_L,\qquad g>0,
\end{equation}
where $P_L$ is orthogonal projection onto $L$. This is a condition on the distribution of executable future tests, not an assumed Fisher matrix. All prefix and suffix noises remain strictly positive. The acquisition order is fixed and every report, preparation and audit is charged as before.

Use the conditional matrices derived by completing squares in the proof of Theorem~\ref{thm:noisy}:
\[
 \Sigma_0=(I_d+\tau^{-2}G^\top G)^{-1},\quad
 C_H=\Sigma_0\tau^{-2}G^\top,\quad
 k=\frac{\Sigma_0a}{\sigma^2+a^\top\Sigma_0a},\quad D=I_d-ka^\top.
\]
Put $A=DC_H$, which is invertible. Choose a fixed orthonormal matrix $J:\R^r\to L$ and define
\[
 c_j=J^\top A e_j,\qquad r_j=\operatorname{rank}[c_1\ \cdots\ c_j],
 \qquad r_0=0.
\]
Then $r_d=r$. These finite matrices are known calibrated coefficients derived from the raw kernels. Computing or storing them is not included in the retained data count.

\begin{theorem}[A rank-sensitive noisy resource law]\label{thm:ranknoisy}
For this experiment and its original binary squared-loss task, with full-history Bayes risk $B_*$,
\begin{align}
 R_\on(M)-B_*&\asymp M^{-2/r},\label{eq:rankonline}\\
 R_\cp(M)-B_*&\asymp M^{-2}.\label{eq:rankcheckpoint}
\end{align}
The constants depend on the fixed raw parameters and frame bound. At each prefix cut $j$ with $r_j>0$ there is an actual dominated continuation measure of exponent $r_j$, with a matching globally implementable cover. The final cut has exponent one. The upper uses only $M$ labels at every cut, not an exact stored prefix followed by dimensional reduction. Its resource count is $d+3$ physical calls, an additional stored phase factor when applicable, and separate program, scratch and numerical precision costs.

Consequently online and checkpoint prediction match in order exactly when $r=1$; for $r>1$ their excess ratio has order $M^{2-2/r}$. An ambient dimension $d>r$ does not change the exponent.
\end{theorem}
\begin{proof}
For $r_j>0$ fix an orthonormal basis matrix $E_j$ of the span of $c_1,\ldots,c_j$ in $\R^r$. The actual task state is
\[
 Z_j=E_j^\top\sum_{i\le j}c_iH_i\in\R^{r_j}.
\]
Zero-rank stages are singletons. Since these spans are nested, its update on the next raw scalar report is
\begin{equation}\label{eq:rankrecursion}
 Z_j=E_j^\top E_{j-1} Z_{j-1}+E_j^\top c_j H_j,
\end{equation}
where the first summand is zero at the initial nonzero rank. The operator multiplying the previous state has norm at most one. Every actual report has all finite moments, also when the coordinates of $H$ are correlated. Lemma~\ref{lem:serialcompander} gives global inward $M$-label quantizers and the candidate Lyapunov bounds for this entire sequence. Applying them after each report retains only the indicated finite index. No exact whitening of $H$ occurs.

At the suffix packet $(V,B)$, put $b'=J^\top B$. The full-history forecast is
\begin{equation}\label{eq:rankresponse}
 P=g_{s_B}\bigl((b')^\top E_dZ_d+B^\top kV\bigr),\qquad
 s_B^2=B^\top\Sigma B,
\end{equation}
where $\Sigma$ and $g_s$ are as in \eqref{eq:posteriornoisy}--\eqref{eq:regresponse}, with the stated $G$ replacement. Formula~\eqref{eq:smoothlower} gives a global Lipschitz constant $1/2$ for this last map in $Z_d$. Quantize its value on $[0,1]$ into $M$ midpoints, and erase the packet and previous state. Thus the final candidate space has exponent one and a bounded cover. The largest exponent among these covers is $r$.

We verify the lower at every active prefix, not only at the last one. For a fixed $j$, restrict the actual vector $(H_1,\ldots,H_j)$ to a compact cube $E$ of positive acquired probability. Conditional on this vector, the remaining coordinates and $V$ have a nondegenerate multivariate Gaussian density: the joint covariance is positive definite because the independent prefix noises have variance $\tau^2>0$ and the suffix noise has variance $\sigma^2>0$. Its conditional covariance is constant, and its mean is linear in the prefix. Hence on any fixed compact cube for $(H_{j+1},\ldots,H_d,V)$ that conditional density has a positive uniform minimum over $E$. Multiplying a uniform probability on this suffix cube by $\lambda$ and including the deterministic complete interface records gives a probability $\eta_j$ and a constant $\alpha_j>0$ with
\[
 \alpha_j\mu_j|_E\otimes\eta_j\le\law(H_{1:j},U_j).
\]
The positive mass $\alpha_j\mu_j(E)$ is retained. Deterministic phases and costs have no additional randomness and are included in the displayed product kernel.

The pushforward of $\mu_j|_E$ by $Z_j$ is dominated by the full Gaussian law of $Z_j$. This is a nondegenerate $r_j$-dimensional Gaussian, since its coefficient matrix has rank $r_j$ and the acquired prefix covariance is positive definite. Its Lebesgue density is bounded above. The pushforward of the restricted measure therefore has small-ball mass at most $C_j u^{r_j}$ at radius $u$, without renormalization.

For fixed complete suffix values the response as a function of the partial state $z$ is
\[
 g_j(z,u)=g_{s_B}\left((b')^\top E_jz+
                  (b')^\top\sum_{i>j}c_iH_i+B^\top kV\right).
\]
It is a bounded Borel extension to all candidate states. On the image of $E$, and on the compact suffix cube, every argument in this formula lies in a common bounded interval. Equation~\eqref{eq:smoothlower} has a strictly positive derivative lower $\gamma_j$ there, also on every segment between two such arguments. Thus
\[
 \|g_j(z,\cdot)-g_j(z',\cdot)\|_{L^2(\eta_j)}^2
 \ge\gamma_j^2\int |B^\top J E_j(z-z')|^2\lambda(dB)
 \ge\gamma_j^2g\|z-z'\|_2^2.
\]
This verifies actual mass and observability with exponent $r_j$. It also shows why a dimension calculation on the hidden state alone would be insufficient.

At the terminal cut, the density argument of Theorem~\ref{thm:noisy} applies for every unit query in $L$: the conditional scalar normal variance is bounded away from zero because $I_d-\Sigma$ is positive definite. Restriction to a fixed bounded normal argument gives an actual forecast submeasure with positive mass and a bounded one-dimensional density. The terminal exponent is therefore one. All the hypotheses of Theorem~\ref{thm:serial} now hold, with $s_*=r$ and $s_T=1$, proving both assertions and the ratio.
\end{proof}

The partial state in \eqref{eq:rankrecursion} is sufficient for the specified response after the remaining reports. It need not determine their entire conditional law: nuisance directions can affect the suffix Gaussian mean while disappearing from the scored response. Accordingly this partial state is not declared to be the full predictive quotient or to supply a free experiment simulator. The full prepared Gaussian experiment still has the report-continuity bound obtained by integrating translated conditional normal densities. A morphism must use its complete transcript and charge whatever extra state its simulator actually retains.

\subsection{A singular acquired law with a noisy continuation}
Fix $d\ge2$, $0<r_j<1/2$ and independent fair digits $\epsilon_{j,n}\in\{0,1\}$. The raw preparation is
\begin{equation}\label{eq:singularraw}
 X_j=1+\sum_{n\ge1}(1-r_j)r_j^{n-1}\epsilon_{j,n},\qquad
 X=(X_1,\ldots,X_d)\in[1,2]^d.
\end{equation}
The compact support of coordinate $j$ is denoted by $C_j$ and its actual preparation law by $\mu_j$. In $d$ paid calls report $X_1,\ldots,X_d$ successively. Then a paid call reports the single packet
\[
 U=(V,B),\qquad V=a^\top X+\sigma Z,\qquad
 B\sim\mathrm{Unif}(\mathbb S^{d-1}),
\]
where $a\ne0$, $\sigma>0$, and $Z,B$ are independent of the preparation. After the forecast, the paid audit reports $Y\sim\mathrm{Bernoulli}(\psi(B^\top X))$. The action class is this fixed acquisition order and all admissible finite-state predictors. Preparation and audit included, the acquisition count is $d+3$.

This is an analog observation kernel: one call reports one real-valued coordinate of the singular preparation. It is not a progressive-digit sensor with infinitely many unpaid reads. Exact Borel processing and finite-precision processing are distinct models. Only a finite label survives each report cut in the theorem below.

Put
\[
 d_j=\frac{\log 2}{\log(1/r_j)},\qquad D_i=\sum_{j\le i}d_j,
 \qquad D=D_d.
\]
These exponents will be proved from the actual cylinder masses and lengths, not postulated as an ambient dimension.

\begin{theorem}[Singular acquisition and terminal geometric expansion]\label{thm:singularserial}
For the raw experiment \eqref{eq:singularraw}, with its binary squared-loss task and full-history Bayes risk $B_*^{\rm sing}$,
\begin{align}
 R_\on(M)-B_*^{\rm sing}&\asymp M^{-2/\max\{D,1\}},\label{eq:singularonline}\\
 R_\cp(M)-B_*^{\rm sing}&\asymp M^{-2}.\label{eq:singularcheckpoint}
\end{align}
Every prefix cut $i$ has a dominated acquired continuation measure with exponent $D_i$ and a recursively usable global cover of that order. The terminal forecast has exponent one. In particular the online and checkpoint laws match when $D\le1$, and strictly separate when $D>1$. The constants retain all fixed ratios, noise and dimension parameters; they are not uniform as a ratio tends to zero or $1/2$.
\end{theorem}
\begin{proof}
We give the covering and mass estimates directly. A level-$n$ cylinder of coordinate $j$ is an interval of length $r_j^n$ carrying mass $2^{-n}$. Distinct level-$n$ intervals have disjoint interiors and gaps at least $(1-2r_j)r_j^{n-1}$. If $0<u<1$, choose $n\ge1$ with $r_j^n\le u<r_j^{n-1}$. An interval of length $2u$ meets at most a fixed number, depending only on $r_j$, of these cylinders. Consequently
\begin{equation}\label{eq:singularball}
 \mu_j([x-u,x+u])\le C(r_j)2^{-n}
                      =C(r_j)(r_j^n)^{d_j}\le C(r_j)u^{d_j}.
\end{equation}
Enlarging the constant covers $u\ge1$ and arbitrary centers. Projecting a Euclidean ball into its coordinates and using the actual product preparation gives a mass bound $C u^{D_i}$ on $C_1\times\cdots\times C_i$.

For a global $M$-point cover of that product, choose
\[
 n_j=\left\lfloor\frac{d_j}{D_i}\log_2 M\right\rfloor,
       \qquad 1\le j\le i.
\]
Retain the endpoint with zero remaining digits in each such cylinder. The number of product representatives is $2^{\sum_j n_j}\le M$, and their coordinate radii obey
\[
 r_j^{n_j}\le r_j^{-1}M^{-1/D_i}.
\]
Thus their Euclidean covering radius is at most $C M^{-1/D_i}$. Strong separation makes the cylinder assignment Borel; nearest-center tie-breaking is another Borel choice. Every representative lies in the same compact product. The raw candidate update appends the next scalar report. It is one-Lipschitz in the preceding coordinates, and with $V_i=1$ its candidate moment bound is automatic. Quantizing after every append therefore yields a genuine recursion with at most $M$ labels at each cut, not a post hoc product code applied to a stored exact vector.

Fix a prefix cut $i$. The remaining coordinates have their unchanged product preparation law. Conditional on the full $x\in[1,2]^d$, the scalar $V$ has a positive Gaussian density, whose minimum on a fixed compact interval $[-T_0,T_0]$ is positive uniformly in $x$. Hence the complete conditional suffix law is bounded below by $\alpha_i\eta_i$, where $\eta_i$ is the product of the remaining coordinate laws, $\mathrm{Unif}[-T_0,T_0]$, the sphere probability and the deterministic interface records. Here $\alpha_i>0$ is the density minimum multiplied by $2T_0$. In particular
\[
 \alpha_i(\mu_1\otimes\cdots\otimes\mu_i)\otimes\eta_i
       \le\law(X_{1:i},U_i).
\]
The mass factor is not removed. On this product the actual conditional response is
$g_i(x_{1:i},u)=\psi(B^\top(x_{1:i},X_{i+1:d}))$; $V$ is observed side information, but cannot improve a conditional forecast once all coordinates have been reported.

All arguments of $\psi$ on the compact candidate product lie in $[-2\sqrt d,2\sqrt d]$. Its derivative there is at least
$\gamma=\tfrac12\operatorname{sech}^2(2\sqrt d)>0$ and at most $1/2$. The mean value theorem and rotational symmetry of the query give
\[
 \|g_i(x,\cdot)-g_i(x',\cdot)\|_{L^2(\eta_i)}^2
 \ge\frac{\gamma^2}{d}\|x-x'\|_2^2.
\]
Together with \eqref{eq:singularball}, this proves the prefix continuation hypotheses with exponent $D_i$. The final update, from the retained full candidate and current packet to its conditional forecast, is globally $1/2$-Lipschitz and takes values in $[0,1]$.

For completeness we establish the terminal lower in the same scored task. Conditional on $X=x$, let $R=\|x\|_2\in[\sqrt d,2\sqrt d]$. A uniform sphere query has scalar projection $T=B^\top x$ with density
\[
 f_R(t)=\frac{c_d}{R}\left(1-\frac{t^2}{R^2}\right)^{(d-3)/2},
 \qquad |t|<R,
\]
where $c_d$ is the positive normalizing constant. This formula follows by slicing spherical surface measure (for $d=2$, by the angular substitution $t=R\cos\theta$). On the fixed interval $|t|\le\sqrt d/2$, these densities are uniformly bounded above and below by positive constants for all admitted $R$. Restricting the actual law to this interval therefore retains positive mass. The change of variable $p=\psi(t)$, whose derivative is bounded below there, gives an actual forecast submeasure with bounded Lebesgue density. The scalar small-ball argument yields $q_M(\law(P))\ge cM^{-2}$; uniform midpoint quantization gives its matching upper. Thus the terminal exponent is one, although every acquired prefix is singular with respect to its ambient Lebesgue measure.

All global and actual-mass hypotheses of Theorem~\ref{thm:serial} have now been verified. Its maximum exponent is $\max(D,1)$, proving \eqref{eq:singularonline}--\eqref{eq:singularcheckpoint}. The new continuous query can increase the last acquired task exponent beyond the preceding singular exponent. Omitting that terminal cut would give the false bound $M^{-2/D}$ when $D<1$.
\end{proof}

The complete suffix kernel is continuous in a changed prefix in total variation with modulus at most $\|a_{1:i}\|_2/(\sqrt{2\pi}\sigma)$: couple the remaining coordinates and query identically, and integrate the translated scalar normal density. This continuity is obtained from the physical noisy channel, not from a nonexistent density of the Cantor preparation. A finite-precision report interface requires its own perturbation certificate; Corollary~\ref{cor:serialerrors} then applies on a declared candidate domain. The finite-bit implementation theorem for progressive sensing remains a separate feasible-resource theorem, not an assertion of unit-cost exact analog arithmetic here.

\section{Adaptive acquisition on nonhomogeneous cylinders}\label{sec:adaptive}
The cut argument concerns a fixed exploration. We now optimize acquisition itself. The raw interface in this section is progressive sensing: a call selects a coordinate and reveals its next unread binary digit. Neither a coordinate value nor an infinite tail is supplied as a pre-decision report. This interface is appropriate to a multiresolution measurement, and its restrictions are part of the theorem. Independent fair unread tails make their future variances depend on depth, not on observed bit values. Thus the policy optimization below is not arbitrary observation-dependent experimental design.

Fix $d\ge1$, positive weights $\omega_j$ with $\sum_j\omega_j=1$, and $0<a\le1/2$. Let
\begin{equation}\label{eq:scales}
 a\le r_{j,l}\le1/2,\qquad
 \ell_{j,0}=1,\quad \ell_{j,k}=\prod_{l=1}^k r_{j,l},\quad
 b_{j,l}=(1-r_{j,l})\ell_{j,l-1}.
\end{equation}
Prepare independent fair bits $\zeta_{j,l}$ and physical coordinates
\begin{equation}\label{eq:cylinders}
 X_j=\sum_{l\ge1}b_{j,l}\zeta_{j,l},\qquad
 \norm{x}_\omega^2=\sum_{j=1}^d\omega_jx_j^2.
\end{equation}
The series converges absolutely and $\sum_l b_{j,l}=1$. After preparation a legal action is either $\mathrm{read}(j)$, which returns the next unread bit of coordinate $j$, or $\mathrm{idle}$, which returns a fixed symbol. Each costs one raw call and one unit of physical acquisition time. The apparatus advances the chosen coordinate's internal read position. That position is part of the physical state, not a tape of past outcomes available to the observer. A policy may choose its actions from its permitted retained history and randomness. It cannot query a digit twice, request the hidden tail in one call, or obtain the terminal audit before predicting.

There are $n$ pre-decision opportunities; shorter acquisition rules are padded with paid idles to this deterministic horizon. A terminal audit returns $X+s\delta\mathbf 1$, with unknown $s\in\{-1,1\}$ and $0\le\delta\le1/8$, only \emph{after} the vector forecast. This is the standard weighted squared reconstruction task. The terminal audit costs one call and does not furnish an input to the prediction algorithm. Preparation also costs one call, so $N=n+2$. Actions and stopping cannot reveal $s$. Decisions may be restricted to the compact box $[-1/8,9/8]^d$, without increasing risk. The oracle risk is zero and the loss is bounded by $25/16$.

A complete history knows a finite prefix in each coordinate. Write $k_j(h)$ for its number of reads and $u_j(h)$ for that prefix. Its conditional mean is
\begin{equation}\label{eq:actualcenter}
 P_j(h)=\tfrac12+\sum_{l\le k_j(h)}b_{j,l}(u_{j,l}(h)-\tfrac12).
\end{equation}
The prefixes and read positions determine every legal future test by direct kernel integration. Conversely, the continuation which reads and scores a chosen coordinate determines its physical conditional law and read position; including the interface gives a predictive realization. For an adaptive policy the nominal acquired score law is exactly
\begin{equation}\label{eq:adaptiveactual}
 \nu_\alpha=\sum_{h}\PP_\alpha(H=h)\delta_{P(h)}.
\end{equation}
The sum is over the finitely many action/report histories of the fixed horizon. Its weights are the products of the policy probabilities and the factors $1/2$ for the actually read digits. Randomizing the policy averages these weights. Prefixes are \emph{not} asserted to be uniform conditional only on the random depths: data-dependent choices may bias that conditioning.

\subsection{A common scale for acquisition and quantization}
Put $s_j(k)=\sqrt{\omega_j}\ell_{j,k}$. Starting from $k^0=0$, increment, at each step, a coordinate with largest $s_j(k_j)$, breaking ties by the smallest index. Set
\begin{equation}\label{eq:greedy}
 R_b=\max_j s_j(k_j^b),\qquad b=0,1,\ldots.
\end{equation}
The total order on merged entries is decreasing numerical scale, then increasing coordinate index, then increasing level; it is compatible with each strictly decreasing coordinate list. The sequence is known from the experiment, not estimated from the observed bits.

\begin{lemma}[Greedy scale and product small balls]\label{lem:greedy}
For every integer $b\ge0$,
\begin{equation}\label{eq:minmaxscale}
 R_b=\min_{\substack{k_j\ge0\ \sum_jk_j\le b}}\max_j s_j(k_j),
 \qquad aR_b\le R_{b+1}\le R_b.
\end{equation}
If $\mu=\law(X)$ and $b_M=\lfloor\log_2M\rfloor$, then
\begin{equation}\label{eq:physicalquant}
 c_{a,d}R_{b_M}^2\le q_M(\mu)\le(d/4)R_{b_M}^2.
\end{equation}
The constants are independent of the weights, the scale sequence and $M$. In particular no limiting dimension is assumed.
\end{lemma}
\begin{proof}
Each list $s_j(0),s_j(1),\ldots$ decreases by a factor in $[a,1/2]$. An increment removes its currently largest remaining entry. The greedy rule therefore removes the $b$ largest entries in the union of the $d$ lists, with a fixed order for ties. Every allocation of $b$ increments removes a prefix of each list. Its maximum remaining entry is at least the $(b+1)$st entry of the merged order; greedy attains it. This proves the first equality, including unused increments by monotonicity. The largest entry removed at the next step is $R_b$, and its successor is at least $aR_b$; this gives the second assertion.

At the greedy allocation $k^b$, every coordinate with $k_j^b>0$ has final weighted cylinder length at least $aR_b$. Indeed its last parent was a maximum at an earlier greedy step, hence was at least the eventual maximum $R_b$. In a fixed coordinate all level-$k_j^b$ cylinders have that same length. They have disjoint interiors because each ratio is at most $1/2$; their endpoints have zero fair-product mass. A ball for $\norm{\cdot}_\omega$ of radius $aR_b/4$ has projection length at most $aR_b/2$ in each weighted coordinate. Here physical length is $\ell_{j,k}$, weighted length is $\sqrt{\omega_j}\ell_{j,k}$, and the ball projection diameter is measured in the latter coordinate. It meets at most three of these intervals in each active coordinate. Each product cylinder has mass $2^{-b}$. Consequently
\begin{equation}\label{eq:productsmallball}
 \sup_x\mu(B(x,aR_b/4))\le3^d2^{-b}.
\end{equation}
Coordinates never split contribute just one interval and do not weaken the bound. The product mass is still $2^{-b}$ because $b$ counts the total number of active binary splits.

Choose $L_d=\lceil\log_2(4\cdot3^d)\rceil$ and $b=b_M+L_d$. Since $M<2^{b_M+1}$, the union of $M$ such balls has mass at most $1/2$. Its complement contributes at least $a^2R_b^2/32$. Iterating \eqref{eq:minmaxscale} gives $R_b\ge a^{L_d}R_{b_M}$, proving the lower with $c_{a,d}=a^{2L_d+2}/32$.

For the upper use the conditional means of the $2^{b_M}$ product cylinders at allocation $k^{b_M}$. Within coordinate $j$ the conditional distribution lies in an interval of length $\ell_{j,k_j^{b_M}}$, so its variance is at most one quarter of that squared length. Summing the weighted variances gives at most $(d/4)R_{b_M}^2$. The centers are physical conditional means; no smooth reference density or Fisher matrix has been introduced.
\end{proof}

\begin{lemma}[Adaptive acquisition variance]\label{lem:acquisition}
For every causal adaptive sensing policy with at most $n$ reads, conditional on a complete realized history the unread digits remain independent and fair. In particular
\begin{equation}\label{eq:tailvariance}
 \frac{\ell_{j,k}^2}{16}\le
 V_j(k):=\frac14\sum_{l>k}b_{j,l}^2
 \le\frac{\ell_{j,k}^2}{4},
\end{equation}
and the full-history acquisition uncertainty satisfies
\begin{equation}\label{eq:adaptivelower}
 \E\norm{X-P(H)}_\omega^2
   =\E\sum_j\omega_j V_j(k_j(H))\ge R_n^2/16.
\end{equation}
This remains true if the sensing controller has unrestricted access to its past reports.
\end{lemma}
\begin{proof}
Let $\mathcal F_t$ be the sigma field generated by the complete action/report history through call $t$ and an independent seed containing all policy coins. Initially, conditional on the seed, the claim is the product preparation. A selected coordinate is a function of revealed data. Its next unread bit is fair and independent of all other unread bits. Revealing it, or performing an idle, preserves the induction. Conditional means and variances are therefore given by \eqref{eq:actualcenter} and the displayed tail series. Its first term has coefficient $b_{j,k+1}\ge\ell_{j,k}/2$, giving the lower in \eqref{eq:tailvariance}. The whole tail has range of length $\ell_{j,k}$, giving the upper. Each realized allocation has sum at most $n$, so \eqref{eq:minmaxscale} bounds its largest weighted length below by $R_n$. Sum the nonnegative variances and then average. Removing the conditioning on coins proves the claim for randomized policies as well.
\end{proof}

\begin{theorem}[Optimized acquisition--memory law]\label{thm:adaptive}
Let $N=n+2\ge2$, $M\ge1$, and $K=\min\{n,\lfloor\log_2M\rfloor\}$. Optimize over all parameter-independent adaptive policies for the raw interface above and all online predictors whose entire data-dependent retained tuple, including any data-dependent acquisition-controller state, has at most $M$ values at every pre-decision cut. Define the analogous checkpoint optimum by allowing the prediction encoder one final access to the complete acquisition history. Then
\begin{equation}\label{eq:adaptivelaw}
 c_{a,d}\{R_K^2+\delta^2\}
 \le\mathcal R_\cp(N,M,\delta)
 \le\mathcal R_\on(N,M,\delta)
 \le C_{a,d}\{R_K^2+\delta^2\}.
\end{equation}
The risk is minimax in the inaccessible sign $s$. A nonadaptive greedy allocation, followed by paid idles, attains the upper while storing only its $K$ observed bits. All constants are uniform in horizon, weights and ratios at fixed $a,d$. The deterministic phase/program is separately charged; it cannot contain observed data.
\end{theorem}
\begin{proof}
Fix any admissible procedure and condition on its independent seed. Averaging the risks at the two signs gives
\[
 \tfrac12\sum_{s=\pm1}\E\norm{X+s\delta\mathbf1-A}_\omega^2
 =\E\norm{X-A}_\omega^2+\delta^2.
\]
Full-history projection and Lemma~\ref{lem:acquisition} make the first term at least $R_n^2/16$. Independently, at the terminal decision it has only $M$ possible vector centers. The law of the physical $X$ is unchanged by any sensing policy or independent seed. Its distortion is therefore at least $q_M(\mu)\ge c_{a,d}R_{b_M}^2$ by Lemma~\ref{lem:greedy}. This is an oracle lower witness, not a substitution of $\mu$ for the actual atomic acquired law \eqref{eq:adaptiveactual}. Since $R_b$ decreases, the larger of $R_n^2$ and $R_{b_M}^2$ is $R_K^2$. These lower bounds hold even for a full-history sensing controller and hence for both classes in the theorem.

For the upper read the $K$ coordinates of the greedy schedule in order. Retain the resulting binary word. At the $t$th read there are $2^t\le2^K\le M$ possible words; all remaining opportunities are paid idles and copy that word exactly. The deterministic phase determines which coordinate each bit belongs to. Output its conditional mean \eqref{eq:actualcenter}. Its nominal squared error is at most $(d/4)R_K^2$. Its unconditional mean error is zero, so each sign adds exactly $\delta^2$. Preparation, all reads and idles, and the terminal audit give exactly $N$ calls. This construction also proves that the upper does not assume a supplied recursively accurate real-valued state.
\end{proof}

For any fixed policy the exact checkpoint identity still separates acquisition and memory:
\begin{equation}\label{eq:policyspecific}
 R_{\cp,\alpha}(M,0)
 =\E\norm{X-P(H)}_\omega^2+q_M(\nu_\alpha).
\end{equation}
In proving \eqref{eq:adaptivelaw} we did not commute an infimum over policies with either summand. Instead, the variance lower holds at every realized allocation and the physical-law converse holds for every encoder.

If all ratios equal $r$ and the weights are $1/d$, greedy depths differ by at most one and $R_K^2$ is comparable, at fixed $d,r$, to $r^{2K/d}/d$. With $r=1/2$ this includes Lebesgue product laws; with $r<1/2$ it includes singular products. Mixed and oscillating sequences are covered by the exact scale sequence even when no single power-law dimension exists. This is a nonuniform geometry statement, not a claim of uniformity as $a\downarrow0$ or $d\to\infty$.

For equal weights $\omega_j=1/d$ and constant ratio $r$, write $K=du+v$ with $0\le v<d$. The greedy allocation has $v$ coordinates at depth $u+1$ and the others at depth $u$, so $R_K^2=r^{2u}/d$. Thus the shorthand $R_K^2\asymp r^{2K/d}/d$ includes a factor between $1$ and $r^{-2}$ caused by the remainder.

\subsection{A common task attaining the deficiency term}
An upper allowance for deficiency cannot force a positive error floor. We give a source--target pair which attains it, while preserving the optimized geometric part of the task.

Add an unknown fixed bit $z\in\{0,1\}$ to the preparation. Exactly one designated paid call, immediately after preparation, reports $z$ with probability $1-2\varepsilon$ and an erasure otherwise, where $0\le\varepsilon\le1/2$. No other pre-decision call depends on $z$. Then allow the same $n$ progressive opportunities. The terminal audit, only after the prediction, scores
\begin{equation}\label{eq:jointtask}
 \ell((X+s\delta\mathbf1,z),(A,c))
 =\tfrac12\norm{X+s\delta\mathbf1-A}_\omega^2+\tfrac12(z-c)^2.
\end{equation}
There are exactly $N=n+3$ raw calls. The exact-bit target changes only the designated report to an exact bit and has the same costs and scored outcome. An admitted simulator must complete this virtual bit report before processing the progressive commands and may not query the terminal audit early.

\begin{theorem}[A matched common acquisition--memory--defect profile]\label{thm:joint}
\textbf{The virtual exact-bit report must be completed before any progressive command or terminal audit.} All comparisons in this theorem impose that order.
For $N\ge3$, $M\ge3$ and $K=\min\{N-3,\lfloor\log_2M\rfloor\}$, the source experiment just defined satisfies
\begin{equation}\label{eq:jointlaw}
 \mathcal R_\cp\asymp_{a,d}\mathcal R_\on
 \asymp_{a,d}R_K^2+\delta^2+\varepsilon,
\end{equation}
where the infima optimize the permitted adaptive sensing and finite-label prediction, and the supremum is over $(s,z)$. Its parameter-uniform causal deficiency to the exact-bit target, with the stated completion interface and common task, is exactly $\varepsilon$. The target-to-source direction is exact by independent erasure. These claims concern this specified family, not every experiment admitting a defect upper bound $\varepsilon$.
\end{theorem}
\begin{proof}
Put the uniform prior on $(s,z)$ for a lower bound on the supremum. Conditional on an erased bit report, all later pre-decision information is independent of $z$, so its Bayes squared error is $1/4$. This contributes $\varepsilon/4$ to the weighted task. The geometric component has the acquisition and physical-law memory lower bounds used in Theorem~\ref{thm:adaptive}; knowing an independent bit or erasure indicator cannot invalidate them. The sign average adds $\delta^2/2$. Taking the larger of the acquisition and memory terms proves the lower.

For the upper retain one of $0,1,\bot$ and $K'=\min\{n,\lfloor\log_2(M/3)\rfloor\}$ greedily acquired digits. The tuple has at most $3\cdot2^{K'}\le M$ values. Predict the observed bit, or $1/2$ on erasure, and use the conditional mean for $X$. The bit contribution is $\varepsilon/4$ for either $z$. Since $0\le K-K'\le2$ and each scale step changes by a factor at least $a$, $R_{K'}^2\le a^{-4}R_K^2$. This proves the upper uniformly in the nuisance parameters and in $N,M$.

To simulate the exact-bit target, keep the observed bit and replace an erasure by an independent fair bit. Its chance of a wrong virtual report is $\varepsilon$ for each $z$. Couple the physical coordinates, all later policy coins and reports until that discrepancy; on its complement the complete common-task experiments coincide. This proves complete-law defect at most $\varepsilon$, uniformly over the admitted target policies. Conversely the two source distributions at the designated completion cut have total variation $1-2\varepsilon$. Their minimax bit-testing error is $\varepsilon$, since the only common report is the erasure. Any simulated exact-bit report would be such a test. Thus its worst-parameter total-variation defect from the exact report is at least $\varepsilon$. The compulsory completion order prevents use of the later audit. Independent erasure of an exact bit implements the reverse direction with zero defect.
\end{proof}

\subsection{A closed finite-bit realization}
The preceding theorems are exact-label statements. Their physical coefficient sequences may be arbitrary real data; such data do not automatically have finite programs. The following result states a finite implementation on a precisely specified region instead of hiding real constants in the machine.

\begin{theorem}[Finite data, scratch, program and time]\label{thm:digital}
The common precision $p$ below is charged for each coefficient and each output coordinate. Bit operations use sequential finite tapes, not unit-cost large integers or an uncharged random-access memory.
For the upper construction in Theorem~\ref{thm:adaptive}, fix its finite greedy schedule of length $K$ and suppose the read-only calibration description supplies dyadic approximations $\widehat b_{j,l}$ for the selected coefficients, with error at most $2^{-p}$, $p\ge1$. The schedule is supplied as certified finite data; no decidability of exact comparison of arbitrary computable reals is assumed. There is a deterministic finite-bit implementation with pre-decision report alphabet $\{0,1,\mathrm{idle}\}$, $K$ retained data bits, and
\begin{align}
 \Gamma&=(K+1)2^{-p},\label{eq:digitalerror}\\
 W&\le C\{p+\log_2(K+d+2)\},\label{eq:workspace}\\
 L_{\rm prog}&\le C(K+d)\{p+\log_2(K+d+N+2)\},\label{eq:program}\\
 T_{\rm bit}&\le C\{N+d(K+d)(p+\log_2(K+d+N+2))\},\label{eq:runtime}
\end{align}
such that its sign-minimax risk is at most
\begin{equation}\label{eq:digitalrisk}
 (d/4)R_K^2+(\delta+\Gamma)^2.
\end{equation}
Here $W$ is additional temporary writable space, $L_{\rm prog}$ includes the schedule and coefficient table, and $T_{\rm bit}$ counts elementary operations on a deterministic sequential-tape implementation. Reads and arithmetic are padded to fixed phase lengths. A supplied clock with $Q\le C(N+T_{\rm bit}+1)$ phases is separately accounted for, giving total persistent cardinality at most $2^KQ$. Physical time is at most $N+T_{\rm bit}$ in the declared units. The terminal scored audit is after the output and is not a pre-decision digital input.

The same construction for Theorem~\ref{thm:joint} uses three additional bit/erasure labels, replaces $K$ by $K'$, and adds the attained bit risk $\varepsilon/4$ with the task weights. Choosing $p\ge\lceil\log_2((K+1)/R_K)\rceil$ puts numerical error at the matched geometric order. These are feasible upper accounts for workspace, program and time, not optimal lower bounds for those resources.
\end{theorem}
\begin{proof}
Append each received bit to a packed retained word and copy that word at idles. The scheduled coordinate and current position are determined by the padded phase, never by unretained observations. At the terminal decision, process coordinates successively. Start the accumulator at $1/2$, scan the retained word and the read-only schedule, and add $\widehat b_{j,l}(\zeta_{j,l}-1/2)$ for the selected coordinate. Dyadic arithmetic with a signed accumulator of $p+\lceil\log_2(K+2)\rceil+O(1)$ bits is exact for these approximate coefficients. Round the output to a $p$-bit dyadic and project to $[0,1]$.

Each coordinate's error from its exact conditional mean is at most $K2^{-p}/2+2^{-p}\le\Gamma$, and so is its weighted norm error. Conditional projection eliminates the cross term with $X-P(H)$. The perturbation and the inaccessible sign together have norm at most $\Gamma+\delta$, proving \eqref{eq:digitalrisk}. The output tape is write-only; later coordinates are computed from the retained word and calibration table, not by rereading previously emitted numbers.

The word uses exactly $K$ data bits. A coefficient and the signed accumulator, with loop indices, fit in \eqref{eq:workspace}. The table of at most $K$ coefficients, coordinate labels and finite loop parameters fits in \eqref{eq:program}. For each of $d$ coordinates perform at most $K+d$ fixed scans/additions, each with the displayed bit length; this deliberately loose sequential bound gives \eqref{eq:runtime}. Heads and loops follow predetermined padded phases, so their positions do not encode extra data. Counting those phases gives the bound on $Q$. A total-state constraint which includes the clock must use $2^KQ$, not just $2^K$. If the model instead supplies a deterministic external clock, its resource coordinate is still recorded. The three-state erasure extension is independent and gives the stated product account.
\end{proof}

\section{The ordinary filtering task}\label{sec:filters}
We verify the block side of Theorem~\ref{thm:transfer} for a structurally different raw experiment. The prediction loss is the usual squared error for a categorical hidden state, not an auxiliary centered probe designed to expose individual coordinates.

Fix $n=d+1\ge2$ and a finite action set. A known matrix $Q_t^a$ has a fixed primitive zero--one support $P$, with $P^m$ entrywise positive, and every nonzero entry of $Q_t^a$ is at least $s_*>0$. Thus every supported product of $m$ such matrices has all entries at least $s_*^m$. Prepare $Z_0$ uniformly on $\{0,\ldots,d\}$. After each transition, state $i$ emits the missing symbol $\bot$ with probability $1-q_i^a$ and otherwise a Gaussian vector of mean $\mu_i^a$ and covariance $\sigma^2I_d$. Require
\begin{equation}\label{eq:gaussparams}
 0<q_-\le q_i^a\le q_+<1,\quad \norm{\mu_i^a}_\infty\le R,
 \quad \sigma>0,
 \quad \abs{\det B_a}\ge b_*>0,
\end{equation}
where the $i$th row of $B_a$, $1\le i\le d$, is $(\mu_i^a-\mu_0^a)^{\mathsf T}/\sigma^2$. The determinant condition is for the acquired lower bound, not for the upper. The finite class parameters and primitive support are fixed; no assertion is uniform as these lower bounds vanish. Time-dependent known emissions can also be used when the same bounds hold.

With respect to Lebesgue measure and a unit atom at $\bot$ the emission density is
\[
 g_i^a(y)=q_i^a(2\pi\sigma^2)^{-d/2}
       e^{-\norm{y-\mu_i^a}^2/(2\sigma^2)},\qquad
 g_i^a(\bot)=1-q_i^a.
\]
For any fixed report-adapted legal exploration, the posterior is positive and has the raw Bayes update
\begin{equation}\label{eq:filterupdate}
 F_t(\pi,a,y)_i=
 \frac{(\pi Q_t^a)_i g_i^a(y)}{\sum_j(\pi Q_t^a)_j g_j^a(y)}.
\end{equation}
The posterior with its declared time/action interface determines all future tests. A terminal audit of $Z_T$, executed after the forecast, separates posterior coordinates, so this is a direct predictive realization. Every transition/report attempt, including a missing report, costs one call. Including preparation and the audit, raw acquisition is $T+2$.

Put $d_H(\pi,\rho)=\osc_i\log(\pi_i/\rho_i)$ and $D_H(\pi)=\osc_i\log\pi_i$. This is a separable metric on the open simplex with its usual Borel structure. The subscripts distinguish spread from the admissibility relation and controller cardinality.

\begin{lemma}[A mass-preserving finite cover]\label{lem:filtercover}
For some $\eta>0$ chosen from the raw class constants below, put $w(\pi)=e^{\eta D_H(\pi)}$. There are at most $M$ Borel code representatives with
\[
 d_H(\pi,Q_M\pi)\le C_{d,\eta}M^{-1/d}w(\pi),\quad
 D_H(Q_M\pi)\le D_H(\pi),\quad (Q_M\pi)_i\ge\pi_i/n.
\]
The uniform seed is a representative. The cover preserves the budget-independent relation
\begin{equation}\label{eq:filterrelation}
 \mathcal A(\pi)=\{\rho:D_H(\rho)\le D_H(\pi),\ \rho_i\ge\pi_i/n\text{ for all }i\}.
\end{equation}
\end{lemma}
\begin{proof}
Choose the least maximal coordinate $j$, and set $v_i=\log(\pi_j/\pi_i)\ge0$. For integer $K\ge1$ define
\[
 \widehat v_i=-\eta^{-1}\log\bigl(\lceil Ke^{-\eta v_i}\rceil/K\bigr),
 \qquad (Q\pi)_i=e^{-\widehat v_i}/\sum_h e^{-\widehat v_h}.
\]
At least one deficit is zero, and there are at most $nK^d$ representatives. The cells are Borel. From $u\le\lceil Ku\rceil/K\le u+K^{-1}$ one obtains
$0\le v_i-\widehat v_i\le e^{\eta v_i}/(\eta K)$.
The oscillation of these errors bounds the projective distance. Inwardness follows from $0\le\widehat v_i\le v_i$. Moreover $\sum_i e^{-v_i}\ge1$ and $\sum_i e^{-\widehat v_i}\le n$, proving the componentwise lower bound. Choose $K=\lfloor(M/n)^{1/d}\rfloor$ when $M\ge n$. For $M<n$ use only the uniform center and $D_H(\pi)\le e^{\eta D_H(\pi)}/(e\eta)$, increasing the fixed constant. The construction and relation are valid on the whole candidate simplex, not only on an observed trajectory.
For completeness, choosing $K=\lfloor(M/n)^{1/d}\rfloor$ when $M\ge n$ gives $K\ge\tfrac12(M/n)^{1/d}$ and $nK^d\le M$, hence $K^{-1}\le2n^{1/d}M^{-1/d}$. This is the stated all-budget conversion.
\end{proof}

\begin{lemma}[Physical block contraction and robust drift]\label{lem:filterblock}
The update \eqref{eq:filterupdate} and relation \eqref{eq:filterrelation} satisfy \eqref{eq:pair}--\eqref{eq:within} with $L=1$, some $\kappa,a<1$ and finite $b,A,B$. Constants depend only on the fixed class parameters and are uniform in deterministic horizon and report-adapted actions. Individual updates need not contract strictly.
\end{lemma}
\begin{proof}
For a nonnegative matrix with no zero column, the derivative of $\log\sum_i e^{z_i}Q_{ij}$ is a convex average of the input direction. Its oscillation cannot increase. Integrating along a segment in log coordinates proves projective nonexpansiveness. Multiplication by a common positive likelihood vector cancels in the log-ratio difference, including at $\bot$.

Choose a fixed cube $G=[-b_1,b_1]^d$. On $G$, uniformly in the hidden state and action, $0<g_-\le g_i^a(y)\le g_+<\infty$. Given every past, $\PP(Y_t\in G\mid\mathcal H_{t-1})\ge p_0=g_-|G|>0$. If $E_j$ denotes the first $j$ reports of a block lying in $G$, then
\[
 \PP(E_j\mid\mathcal H_s)
 =\E[\one_{E_{j-1}}\PP(Y_{s+j}\in G\mid\mathcal H_{s+j-1})\mid\mathcal H_s]
 \ge p_0\PP(E_{j-1}\mid\mathcal H_s).
\]
Thus the good block $E_m$ has conditional mass at least $p=p_0^m$, without an independence assumption. Its unnormalized product matrix has entries in $[l,u]$, where $l=s_*^m g_-^m$ and $u=g_+^m$. Every column derivative distribution then dominates $(l/u)$ times the same input probability distribution. Removing that common mass gives projective gain $\chi=1-l/u<1$. Off the block use nonexpansiveness, so $\kappa=1-p(1-\chi^2)<1$.

For any adapted relation-preserving perturbed sequence, the good block gives componentwise
$z_j\ge n^{-1}(g_-/g_+)z_{j-1}Q_{s+j}^{a_{s+j}}$.
Its final coordinates are at least $c_*=n^{-m}(g_-/g_+)^ms_*^m$. Therefore $D_H(z_m)\le R_*:=\log(1/c_*)$. This proves why a cone-mass constraint, not spread inwardness alone, is needed through a sparse product.

For an arbitrary report let $V_a(y)=\osc_i\log g_i^a(y)$ and $C_s=\log(n/s_*)$. Then $D_H(F(z,a,y))\le D_H(z)+C_s+V_a(y)$, and admissible rounding does not increase spread. Gaussian log-likelihood ratios are affine functions of $y$, and the missing ratios are bounded. Hence for any fixed $\eta_0>0$,
$\mathcal M:=\sup_{x,a}\int e^{2\eta_0 V_a(y)}J_a(dy\mid x)<\infty$.
For example $V_a(y)\le v_0+v_1\norm{y}_1$, and
$\E e^{c\norm{Y}_1}\le2^d e^{cdR+dc^2\sigma^2/2}$ gives a uniform finite bound on the continuous component. Conditional iteration yields $\E e^{2\eta_0 Z}\le C_0$ for the block sum $Z=\sum_j(C_s+V_{a_{s+j}}(Y_{s+j}))$.

Choose $0<\eta\le\eta_0$ so that $(\eta/\eta_0)(C_0-1)\le p/2$. Convexity in the exponent gives
$e^{2\eta Z}-1\le(\eta/\eta_0)(e^{2\eta_0Z}-1)$.
Consequently $\E[e^{2\eta Z}\one_{E_m^c}\mid\mathcal H_s]\le1-p/2$. On $E_m$ the envelope is at most $e^{\eta R_*}$; off it, it is at most $w(z_0)e^{\eta Z}$. This proves \eqref{eq:drift} with $a=1-p/2$, $b=e^{2\eta R_*}$. The same growth estimate for every prefix gives \eqref{eq:within} with $A=C_0$, $B=0$. The proof applies to every adapted admissible perturbation under the true report law, not just to the chosen encoder.
\end{proof}

\begin{theorem}[Sharp state order for squared filtering loss]\label{thm:brier}
Use the ordinary Euclidean norm inherited from $\R^{d+1}$ on the simplex. Euclidean projection of any forecast onto that closed convex simplex cannot increase its squared loss against a categorical outcome.
Under the preceding raw assumptions, fix $T\ge m$ and any report-adapted exploration independent of the predictor. Let
\[
 B_T^{\rm fil}=\E\norm{e_{Z_T}-\pi_T}_2^2.
\]
For prediction of the categorical state under loss $\norm{e_{Z_T}-A}^2_2$, with $A$ in the closed simplex,
\begin{equation}\label{eq:brier}
 cM^{-2/d}\le R_\cp(T,M)-B_T^{\rm fil}
 \le R_\on(T,M)-B_T^{\rm fil}\le CM^{-2/d}.
\end{equation}
The constants are uniform in $T$ and the specified exploration at fixed dimension, support, primitive length and parameter bounds. The readout state is a prediction label, not an unrestricted optimized controller. Controller and deterministic clock resources remain separate.
\end{theorem}
\begin{proof}
Differentiating the softmax along a log-coordinate segment gives derivative components $\pi_i(h_i-\sum_j\pi_jh_j)$. The sum of their absolute values, and hence their Euclidean norm, is at most $\osc(h)$. The posterior readout is therefore $1$-Lipschitz from $d_H$ to Euclidean norm. Lemmas~\ref{lem:filtercover} and \ref{lem:filterblock}, with the exact uniform seed, verify every upper hypothesis of Theorem~\ref{thm:transfer} and give $CM^{-2/d}$.

For the lower, retain only the actual event that the last $m$ reports lie in $G$. It has mass at least $p$. Conditional on the prefix and the first $m-1$ reports of this event, the predictive weights $r=\pi_{T-1}Q_T^{a_T}$ have coordinates at least $c_r=(g_-/g_+)^{m-1}s_*^m$. The last action is selected \emph{before} the final report. Its log posterior ratios are $B_{a_T}y$ plus a fixed vector, so the map from $y$ to the first $d$ posterior coordinates is injective. The softmax derivative in those coordinates is $\operatorname{diag}(\pi_1,\ldots,\pi_d)-vv^{\mathsf T}$, where $v=(\pi_1,\ldots,\pi_d)$. The matrix determinant lemma gives determinant $\prod_{i=0}^d\pi_i$. Thus the report-to-posterior Jacobian has absolute determinant at least $b_*c_p^{d+1}$ on $G$, where $c_p=c_r g_-/g_+$.

The last conditional report density is at most $g_+$. The resulting posterior submeasure has density at most $g_+/(b_*c_p^{d+1})$ in its first $d$ coordinates. Integrating over the conditioning prefix with its event indicator preserves this bound. The submeasure has its actual mass at least $p$; it is not normalized to mass one. An ambient Euclidean ball projects to a $d$-ball of the same radius, so its mass is at most $C_1r^d$. The union-of-balls lower in Theorem~\ref{thm:transfer} gives $cM^{-2/d}$. Projection for the categorical task identifies the common baseline and completes the proof.
\end{proof}

Report continuity is also explicit: mixtures of the same emission kernels have total variation at most $\frac12\norm{\pi Q-\rho Q}_1\le\frac12d_H(\pi,\rho)$. The terminal categorical audit obeys the same bound without $Q$. The total stopped-law comparison in Section~\ref{sec:composition} therefore applies. This risk theorem does not require a preassigned Fisher geometry; the exponent comes from a positive submeasure of the actual posterior law.

For a finite numerical interface, a bounded input box is essential. If all continuous reports before $T$ are clipped outside $[-U,U]^d$, then the failure probability is at most $2dT\exp(-(U-R)^2/(2\sigma^2))$. Choosing
\begin{equation}\label{eq:clipping}
 U=R+\sigma\sqrt{2\log(2dT/\zeta)}
\end{equation}
limits this defect to $\zeta\in(0,1)$. In particular $\zeta=M^{-2/d}$, with a harmless fixed reduction when $M=1$, requires growth of order $\sqrt{\log T+\log M}$ at fixed parameters. If intervals $[l_i,u_i]$ enclose unnormalized posterior log weights with width at most $h$, deficits $d_i=\max(0,\max_jl_j-u_i)$ satisfy $0\le d_i\le v_i$ and $v_i-d_i\le2h$. Applying the compander to these lower deficits preserves \eqref{eq:filterrelation} and gives local allowance $u=2h$. Proposition~\ref{prop:repair} charges clipping by bounded loss. These are enclosure and defect certificates; unlike Theorem~\ref{thm:digital}, they do not assert a complete optimal Gaussian bit-complexity region.

The density in the lower proof is with respect to the chart $(\pi_1,\ldots,\pi_d)$, with $\pi_{d+1}=1-\sum_{i\le d}\pi_i$; its Euclidean metric is comparable to the full simplex metric. Choosing the clipping failure probability $\zeta=O(M^{-2/d})$ adds at most the bounded-loss constant times $\zeta$ and therefore preserves the displayed excess order.

\section{Stopped interfaces and comparison of experiments}\label{sec:composition}
The lower and upper sides must refer to the same task under any experimental reduction. A complete-law defect is stronger than a predictive discrepancy and is added after root metric errors have been combined.

\subsection{Bounded stopped prediction}
For a visible stopping time $\tau\le mK_0$, define the stopped sigma field $\mathcal H_\tau$, a declared bounded task $Y^\tau$, and
\[
 P_\tau=\E[Y^\tau\mid\mathcal H_\tau],\qquad
 B_\tau=\E\norm{Y^\tau-P_\tau}^2.
\]
At $\tau=0$ these mean the actual preparation and its legal terminal task. Assume $P_\tau=f_\tau(S_\tau)$, with a common Lipschitz constant, and that every stopped exploration is the prefix of a specified legal continuation to $mK_0$ for which the block certificate holds. This continuation is used only for estimates; post-stop calls and information are not used by the actual predictor.

\begin{proposition}[A selection-sensitive stopped bound]\label{prop:stopped}
Under the upper hypotheses of Theorem~\ref{thm:transfer}, put $\rho=r_M+u$ and $K_\tau=\lceil\tau/m\rceil$. There are explicit finite constants $c,C$, depending only on the block certificate, such that with
$V_0=d(S_0,\widehat S_0)^2+c\rho^2w(\widehat S_0)^2$,
\begin{equation}\label{eq:stoppedbound}
 \E d(S_\tau,\widehat S_\tau)^2
 \le C\{\E V_0+\rho^2\E K_\tau\}.
\end{equation}
Consequently $e_0=O(\rho)$ and $\E K_\tau\le L_0$ give
$R\le B_\tau+(C'\rho\sqrt{1+L_0}+v)^2$.
If the terminal decoder has a prediction label of size $M$ and a retained stopping clock of size $Q$, its checkpoint identity uses $q_{MQ}(\law(P_\tau))$ when this is the whole accessible tuple. A free declared interface instead requires conditioning on that interface.
\end{proposition}
\begin{proof}
Here are constants and the optional-sum argument, to distinguish the claim from substituting a random time into a deterministic estimate. Put
\[
 \bar\kappa=(1+\kappa)/2,\quad C_Y=\bar\kappa/(\bar\kappa-\kappa),
 \quad\gamma=\max\{\bar\kappa,(1+a)/2\},
\]
\[
 c=\max\{1,C_Y\Lambda_m^2A/(\gamma-a)\},\qquad
 c_0=C_Y\Lambda_m^2(B+1)+cb.
\]
The conditional shadow bound, $r_M\sqrt{x}+u\le\rho\sqrt{x+1}$ and
$(\sqrt\kappa d+z)^2\le\bar\kappa d^2+C_Yz^2$ give
\[
 \E[V_{k+1}\mid\mathcal H_{km}]\le\gamma V_k+c_0\rho^2,
 \quad V_k=d(S_{km},\widehat S_{km})^2+c\rho^2w(\widehat S_{km})^2.
\]
Multiply the rearranged inequality by the predictable indicator $\one_{\{k<K_\tau\}}$ and sum the finitely many terms. This gives
$(1-\gamma)\E\sum_{k<K_\tau}V_k\le\E V_0+c_0\rho^2\E K_\tau$.
For a raw stopping time, the within-block telescoping bound and the inequality that a maximum of nonnegative terms is at most their sum give
\[
 \E[\max_{1\le j\le m}d(S_{km+j},\widehat S_{km+j})^2\mid\mathcal H_{km}]
 \le C_1V_k+C_2\rho^2,
\]
where $C_1=2m(L^{2m}+\Lambda_m^2A/c)$ and $C_2=2m\Lambda_m^2(B+1)$. Sum only over entered blocks; their indicator is known at each block start. At $\tau=0$ use $V_0$. The resulting bound is
\[
 (1+C_1/(1-\gamma))\E V_0+
 (C_1c_0/(1-\gamma)+C_2)\rho^2\E K_\tau,
\]
which proves \eqref{eq:stoppedbound}. Stopped conditional projection proves the risk statement. The information-cut identity follows by counting every available terminal state before applying Lemma~\ref{lem:cut}; it is not an unqualified $M$-center formula when the clock carries information.
\end{proof}

\subsection{Measurable report coupling}
\begin{lemma}[Measurable maximal report coupling]\label{lem:coupling}
For probability kernels $P_z,Q_z$ on a standard Borel report space, there is a jointly measurable coupling whose disagreement probability is $\norm{P_z-Q_z}_{\TV}$ for every parameter $z$.
\end{lemma}
\begin{proof}
One construction uses nested finite partitions generated by a countable separating algebra. Relative to $R_z=(P_z+Q_z)/2$, their atomwise density ratios are bounded measurable functions. Martingale convergence gives a jointly measurable clipped limsup version $p=dP_z/dR_z$ for each $z$; put $q=2-p$. The common kernel $\mu_z=\min(p,q)R_z$ has mass $c_z$. The coupling
\[
 \operatorname{diag}_*\mu_z+
 \frac{(P_z-\mu_z)\otimes(Q_z-\mu_z)}{1-c_z}\quad(c_z<1)
\]
and the diagonal law for $c_z=1$ are measurable kernels. The residual measures are mutually singular, so their product has zero diagonal mass. The disagreement probability is $1-c_z=\TV(P_z,Q_z)$, with $\TV=\sup_A|P(A)-Q(A)|$.
\end{proof}

Suppose report kernels satisfy $\TV(J(\cdot\mid x,a),J(\cdot\mid z,a))\le L_Jd(x,z)$, and task kernels have constant $L_Y$. A preparation mismatch is either absent or explicitly at most $\varepsilon_0$. Successive maximal couplings, driven until first disagreement by the same declared controller and stopping rule, give complete stopped report-and-task defect at most
\begin{equation}\label{eq:reporttv}
 \min\{1,\varepsilon_0+L_J\E\sum_{t=1}^\tau d(S_{t-1},\widehat S_{t-1})
                         +L_Y\E d(S_\tau,\widehat S_\tau)\}.
\end{equation}
Indeed before disagreement both transcripts, actions and stopping decisions coincide. Continue a physical-law shadow after disagreement and sum first-disagreement probabilities, dropping the agreement indicator. This preserves the unconditional shadow estimates. A stopped sum is not replaced by $\E\tau$ times an unconditional deterministic-time bound. A finite complete-law bound is not infinite-path or large-deviation equivalence.

\subsection{Resource-certified morphisms}
For clarity we state all components of the comparison used here. A source experiment $\mathfrak E$ implements a target $\mathfrak F$ on a common parameter set by parameter-independent transducer initialization, state-update, source-action and report-transformation kernels. These kernels induce a strategy lift from the declared target class to legal source strategies. Each virtual completion time $\sigma_i$ is a source stopping time bounded by the given physical budget; for every permitted target stop $\tau$, $\sigma_\tau$ is a legal source stop and no later source information enters its decision. Legal actions and the available interface agree after each transformed prefix. Virtual costs, preparation and failures are included.

There is one Borel common scored-outcome map and one decision space, with loss in $[0,L_*]$. Uniformly over the common parameters, declared target policies and permitted bounded stops, the joint law of virtual history, permitted public policy coins and scored outcome differs from its target law by at most $\varepsilon$ in total variation. A transducer relying on a real hidden history or unspecified phase is not a finite implementation. Let its state count be $R$, target prediction count $M$, controller count $C_{\rm ctr}$ and phase/clock count $Q$. Workspace, read-only program, precision, raw calls and physical time are separate coordinates. Submitted Supplement S~\cite{companion} gives the fully typed standard Borel transcript construction; these requirements are precisely the ones needed in the following implication.

\begin{theorem}[Common-risk composition and reverse transfer]\label{thm:morphism}
Such a certificate implements a target risk $r_\lambda$ by a serial source tuple of at most $MR C_{\rm ctr}Q$ states and risk at most $r_\lambda+L_*\varepsilon$. Defects of two composable certificates add, capped at one, provided the lifted policies, stopping rules and common-task maps lie in the next certified class. Serial state overheads multiply; raw costs sum over executed subroutines, and nested per-call upper bounds may multiply.

If every source $K$-state procedure can conversely be implemented in the target with at most $\Psi(K)$ states and complete defect $\varepsilon_-$, then for the same task and risk criterion
\begin{equation}\label{eq:morphismlower}
 R_E(K)\ge R_F(\Psi(K))-L_*\varepsilon_-.
\end{equation}
In particular every acquired continuation lower of Lemma~\ref{lem:cut} transfers with that inflated budget and the target baseline left unchanged. A block upper transfers as
\begin{equation}\label{eq:morphismupper}
 R_E(MRC_{\rm ctr}Q)
 \le B_T+(L_fE_{M,u}+v)^2+L_*\varepsilon.
\end{equation}
The same assertions hold for minimax risks when both certificates are uniform over the same parameter set.
\end{theorem}
\begin{proof}
Run the physical subroutine and retain the tuple of predictor, simulator, controller and clock states. The action and stopping conditions make the procedure executable; its cardinality is the displayed product. Applying the common decision kernel preserves total-variation closeness. The bounded-loss inequality for our convention gives a difference at most $L_*\varepsilon$.

For composition insert the intermediate stopped joint law and use the triangle inequality and contraction of total variation under the downstream parameter-independent transformation. Class closure is needed here for adaptive lifted policies. Retaining both serial tuples proves the storage upper. All actual calls, waiting times and failures are summed; shared program or workspace can sometimes be reused, so additive serial descriptions are only upper accounts.

For the reverse statement apply the upper implication to every source procedure: $R_F(\Psi(K))\le r_E+L_*\varepsilon_-$. Infimize over source procedures. Uniformity permits the same argument after taking the parameter supremum. No separately recomputed source and target Bayes baselines are subtracted. This proves both \eqref{eq:morphismlower} and \eqref{eq:morphismupper} in their stated directions.
\end{proof}

The exact-bit and erasure pair in Theorem~\ref{thm:joint} verifies this definition with a prescribed completion clock, parameter-independent bit transformation and a common terminal audit; it is not just a marginal report comparison. The Gaussian report comparison verifies \eqref{eq:reporttv} for the same declared controller. These are distinct uses: the latter does not manufacture a reverse certificate for all Gaussian source algorithms. Likewise a product $RC_{\rm ctr}Q$ is not necessary merely because one implementation stores it. Only an executable challenge separating those registers, or an appropriate reverse theorem, provides that lower bound.

The target class in the exact-bit comparison is all progressive policies with the compulsory early completion order of Theorem~\ref{thm:joint}. In the Gaussian filtering comparison it is the fixed report-adapted physical controller of Theorem~\ref{thm:brier}, with the bounded visible stopping class of Proposition~\ref{prop:stopped} when stopping is used. In Corollary~\ref{cor:noisynumeric} it is the fixed-order regression acquisition and its finite-state prediction class. No implication identifies these different classes.

\section{Antecedents, distinctions and quantitative scope}\label{sec:literature}
\subsection{What the transfer theorem adds}
Conditional projection, small-ball quantization and finite-center approximation are classical; see Graf and Luschgy~\cite{graf}. Predictive representations and future equivalence are developed by Littman, Sutton and Singh~\cite{psr} and Shalizi and Crutchfield~\cite{shalizi}; measurable kernel sufficiency has a categorical treatment in Fritz~\cite{fritz}. Blackwell~\cite{blackwell}, Le Cam~\cite{lecam} and Torgersen~\cite{torgersen} provide the statistical-comparison background. Proposition~\ref{prop:quotient} makes explicit a measurable-image certificate and common report versions; it does not originate predictive sufficiency. The classical total-variation inequality is not the new part of Theorem~\ref{thm:morphism}.

Average contraction of random maps is also classical~\cite{diaconis}. The constructive part of Theorem~\ref{thm:transfer} uses a stronger, specifically physical-law certificate: a cover-preserved relation must control every adapted candidate perturbation across a block. The positive-cone lower bound in Lemma~\ref{lem:filtercover} shows how such a relation can be manufactured for primitive sparse products. Its role is to preserve accumulated positivity while rounding after each raw call. It is sufficient, not an intrinsic necessary condition. Proposition~\ref{prop:repair} allows flagged implementation failures, but does not show that arbitrary local covers or ordinary inwardness imply robust block drift.

The opposite direction now uses actual distributions of continuation-response functions under a joint-law domination certificate. Unlike exact tag separation, it allows a continuum of noisy response functions and charges distortion in the same terminal squared task. Unlike a last-marginal small-ball estimate, it can distinguish a large information cut from an atomic terminal answer. The proof is a Hilbert-space quantization argument; the random-access example explicitly uses an established entropy mechanism~\cite{nayak}. The companion's earlier exact-tag and soft-information arguments remain complementary, rather than being renamed as new theorems.

The optimized acquisition result uses a different constructive mechanism: exact refinement of acquired prefixes, with no rounding of an unretained real posterior. Quantization on Moran sets and measures has an extensive independent theory; for example Zhu~\cite{zhu-moran} proves comparability of optimal Voronoi-cell errors under doubling and separation assumptions. Lemma~\ref{lem:greedy} is a self-contained nonasymptotic product estimate. Its use in Theorems~\ref{thm:adaptive} and \ref{thm:joint} is to identify one scale which also controls every adaptive allocation of paid digit queries. We do not claim a general new quantization theorem for all Moran measures, nonproduct sources or arbitrary adaptive partitions.

\subsection{Decoder side information and functional source coding}
In the Wyner--Ziv problem~\cite{wynerziv}, the encoder observes a memoryless source block $X^n$, while the decoder additionally observes $Y^n$; the rate--distortion expression minimizes $I(X;W)-I(Y;W)$ over a Markov auxiliary $W-X-Y$ and a decoder reproduction satisfying the distortion constraint. At our cut, $H$ is the encoder observation, $U$ is decoder side information, and $z(H,U)$ is the remote conditional-mean task. These identifications are substantive: fixing a causal cut produces an ordinary side-information coding problem. Formula~\eqref{eq:exactfunctional} is its direct finite-code variational form, not a new source-coding principle. Unlike an asymptotic rate, $M$ bounds the entire retained tuple in one actual experiment, and the task and its Bayes baseline remain fixed through the cut.

Orlitsky and Roche~\cite{orlitsky} study reliable computation of a function of the sender's and receiver's combined data. Their problem is closer than ordinary reconstruction of the sender's source: future reports change the function that the earlier information must support. Our response function $z(h,\cdot)$ records exactly that dependence. The difference is not the observation that functions can be compressed. It is the quantitative use of an actual product sublaw to put all responses in a common $L^2(\eta)$ metric, together with retained submeasure mass and recursively attainable bounds under the stated raw acquisition interface. Zero-error distinguishability and asymptotic functional coding are not being asserted equivalent to squared-distortion Hilbert quantization.

Nor is finite blocklength, by itself, a novelty claim. Watanabe, Kuzuoka and Tan~\cite{watanabe} give nonasymptotic achievability and second-order bounds for side-information problems, including Wyner--Ziv coding. Those bounds concern coding probabilities and rates under their information-spectrum formulation. Here \eqref{eq:cutsandwich} instead compares a hard one-cut cardinality directly with a finite-measure geometric distortion, and \eqref{eq:conditionalcut} localizes when the global reference product fails. Lemma~\ref{lem:minorization} derives the reference from physical report kernels. It does not improve the general finite-block coding bounds of that literature.

The recursively executable conclusion is separate from all these one-cut reductions. Theorem~\ref{thm:chart} explicitly stores quantized incoming coordinates, proves an all-tail moment bound, and gives equivalence of online risk and continuation width on its declared correlated sequential class. Theorem~\ref{thm:noisy} computes a positive-noise, continuous-query example in which the same task has different online and checkpoint exponents. Gaussian conditioning, companding and density quantization are standard tools. The claimed synthesis is the raw-kernel continuation lower together with the legal intermediate-state upper and their common-task separation; no independent priority certification is inferred from this comparison.

\subsection{Allocation and successive refinement}
Incremental allocation for a separable discrete resource constraint has classical marginal-analysis antecedents, including Fox~\cite{fox}. Our merged-list argument is an elementary exchange argument in that tradition. In particular, it does not introduce greedy bit allocation for independent components. The objective in \eqref{eq:minmaxscale} is the largest remaining weighted scale, not an unspecified average-rate Lagrangian. Its value also controls the physical-law small balls and the residual uncertainty under every paid progressive policy. These additional links are what make one scale govern both $N$ and $M$ in \eqref{eq:adaptivelaw}.

Equitz and Cover~\cite{equitz} characterize successive refinement by compatible rate--distortion-optimal reproductions, with the coarse reproduction conditionally generated from the finer one. Their staged code describes a known source block at asymptotic rates. Here each new digit must first be physically acquired, and all already acquired digits count against a hard state cardinality rather than an entropy-coded average rate. Exact prefixes provide an executable refinement mechanism, but we do not deduce that arbitrary sources are successively refinable. Nonhomogeneous coordinate scales and singular factors remove a stationary high-rate formula, not the independence assumption. Observed values cannot change future scale values in this model; an extension to genuinely value-dependent experimental design requires another acquisition converse.

\subsection{Operationally close work}
Recursive nonlinear-filter quantization predates this work. Pag\`es and Pham~\cite{pages} and Sellami~\cite{sellami} transport quantization error through filtering recursions. Their approximation constructions do not by themselves identify a fixed cardinality for every retained predictive label or a matching actual-posterior converse. Here Lemmas~\ref{lem:filtercover} and \ref{lem:filterblock} provide that count, and Theorem~\ref{thm:brier} gives the ordinary filtering excess. No novelty is claimed for the Bayes formula, softmax derivative or projective nonexpansiveness separately.

Theorem~9 of Subramanian, Sinha, Seraj and Mahajan~\cite{ais} bounds approximate dynamic-program values from an approximate information-state generator and reward/transition discrepancies. Its policy-loss conclusion concerns control optimization. Theorem~16 of Kara and Y\"uksel~\cite{kara} derives finite-window control bounds under their stability, finite-alphabet and discount hypotheses. Our same-controller physical-law estimate does not imply either optimal-control conclusion. Conversely their conditions do not supply the actual acquired-law lower bound or the hard predictive-state alphabet of \eqref{eq:brier}.

Reachable-belief covering is not a new complexity viewpoint. Zhang, Littman and Chen~\cite{zhang2012} connect covering numbers to POMDP planning and no-reset learning. Theorem~1 of Zhang, Hsu and Lee~\cite{zhang2014} gives a heuristic-guided planning-time upper of order $h\,C(\delta/2)^2$, with depth and resolution chosen from the discount and target value error. Its cover concerns a heuristically reachable search space, not the probability mass actually acquired under a fixed experiment. Grover and Dimitrakakis~\cite{grover} provide adaptive belief discretization and explicit planner-memory/value-error bounds. They therefore cannot be dismissed as not addressing memory; the distinction is planning storage and value approximation versus per-cut predictive alphabet, acquisition uncertainty and a matching geometric lower.

Recent work likewise uses belief geometry operationally. Zhu and Lu~\cite{zhulu} develop offline POMDP evaluation bounds through belief-space coverage and policy stability. This is a learned off-policy evaluation problem, whereas \eqref{eq:central} fixes known kernels and \eqref{eq:adaptivelaw} optimizes a specified progressive sensor. Anjarlekar, Etesami and Srikant~\cite{anjarlekar} use finite-history superstates and policy-based learning with approximation controlled by history length. We do not obtain their learning guarantees from our known-kernel digital implementation. These comparisons are to their published conference versions, not predictions of future results.

Finite-state controller synthesis also has an independent objective. Andriushchenko, \v{C}e\v{s}ka, Junges and Katoen~\cite{inductive} synthesize controllers for POMDP specifications. Hud\'ak and collaborators~\cite{hudak} extract finite-state controllers from learned recurrent policies and evaluate them, including hidden-model robustness. These are substantial finite-controller results, not instances of a terminal codebook. Our converse fixes a common statistical task and the information crossing a cut; it does not establish controller synthesis completeness, robustness for unknown dynamics, or superiority in computational time.

Zero-delay coding and nonanticipative rate-distortion theory address related causal constraints. Wood, Linder and Y\"uksel~\cite{wood}, Ghomi, Linder and Y\"uksel~\cite{ghomi}, and Charalambous, Stavrou and Ahmed~\cite{nonanticipative} study structural or information-constrained formulations. A mutual-information budget is not automatically a hard retained alphabet; nor does our finite-horizon state-order result prove existence of optimal stationary average-cost codes.

\subsection{Scope of the resulting laws}
The new obstruction and the sufficient construction deliberately have different assumptions. Global joint-law domination may have very small mass or fail entirely. Conditional cuts can remain useful, as Example~\ref{ex:diagonal} proves. Theorem~\ref{thm:chart} proves quantitative completeness on its explicit sequential-chart class; passing every available cut test is not proved sufficient outside such a constructively verified class. The profile \eqref{eq:adaptivelaw} is optimized over every adaptive allocation in its raw sensing class, but not over arbitrary interventions, unknown-kernel learning or unbounded stopping. It includes only the stated deterministic phase convention; total persistent state includes any separately stored clock.

For the Gaussian realization, constants deteriorate with $1-\sqrt\kappa$, $1-a$, primitive length, positivity floors, acquired mass and the inverse determinant in \eqref{eq:gaussparams}. For the multiscale realization they depend on $a,d$, not on a nonexistent assumed limiting dimension. Theorems~\ref{thm:digital} and \ref{thm:morphism} charge workspace, program, simulator and time, but do not prove the full optimal feasible region for those coordinates. A lower bound on an output grid is not a lower bound on each internal arithmetic operation. A calibration or deficiency \emph{allowance} alone is not a floor; \eqref{eq:jointlaw} attains such terms only in its explicit inaccessible-sign and erasure family.

Thus the central connection is mathematical rather than terminological: the prepared kernels determine future tests, their measurable predictive realization and actual acquired laws; compatible finite recursions construct the upper; dominated continuation cuts and acquisition variance supply converses; and typed causal morphisms transfer both at their declared resources. The class equivalence and the noisy separation are proved without an independence or contraction assumption on acquired prefix coordinates, but do not constitute a necessary-and-sufficient classification of every recursive acquired geometry.

\paragraph{Serial acquired exponents.}
Theorem~\ref{thm:serial} does not claim new small-ball quantization, Minkowski inequalities or linear error propagation. Its conclusion concerns the same scored experiment at every retained-state cut: observable acquired submeasures supply a lower of each exponent, while global inward covers and a candidate Lyapunov envelope supply a single causal construction. This establishes a conditional compatibility criterion and an explicit ratio of online to checkpoint excess. Neither an asymptotic decoder-side-information rate nor a separate optimal quantizer at each time alone supplies that recursion. The rank-sensitive example verifies the entire sequence of cuts, not only a terminal covariance. In the singular example the preparation geometry and the later task geometry have different exponents, so the maximum must include the final cut. The scope of the criterion is the stated raw-verifiable certificate class; it is not an unconditional duality theorem for all experiments.

\appendix
\section{Retained singular dynamics and interface conventions}\label{sec:appendix}
\subsection{The nonreset singular realization}
Submitted Supplement S~\cite{companion}, Section~5, supplied in full and intended to be refereed with this article, supplies a second direct verification of the block hypotheses, distinct from Gaussian filtering. We record its raw data and the exact point of attachment so that its status is not confused with the adaptive progressive-sensor class of Section~\ref{sec:adaptive}.

Take ratios $r_l\in[1/5,1/3]$, $\ell_l=\prod_{j\le l}r_j$ and a fair binary Cantor coordinate $x=\sum_l(1-r_l)\ell_{l-1}\zeta_l$. Complete its domain by adjoining every finite-prefix fair-tail mean. A prepared mode $I\in\{0,1\}$ is reported, followed by $B$ paid initial prefix reads. At a later opportunity an independent active indicator has any deterministic probability. Inactivity is an exact copy but costs a call. On activity, mode zero either prepends two fair reported bits, deletes the first bit within mode zero, or changes to mode one while deleting; the probabilities are $(1-\alpha_t)(1-\gamma_0)$, $(1-\alpha_t)\gamma_0$, and $\alpha_t$, with $\gamma_0=1/4096$. Mode one prepends six reported fair bits and moves to zero with probability $1/2$, or copies within mode one with probability $1/2$. Deletion does not reveal the deleted bit; both mode changes transport a nonconstant old tail.

Within each mode give the completed coordinate Euclidean distance weighted by $v_0=1800$, $v_1=1$; between modes use distance $2\sqrt{1800}$. The relation is the full same-interface candidate domain and $w=1$. On the completed domain deletion has gain at most $30$ and prepending $s$ fixed bits has gain at most $6\,3^{-s}$. The matrix of physical expected squared gains is
\[
 C_t=\begin{pmatrix}(1-\alpha_t)85/128&900\alpha_t\\2/59049&1/2\end{pmatrix},
 \qquad C_t\binom{1800}{1}\le\frac{17}{20}\binom{1800}{1}.
\]
Cross-mode pairs have squared gain at most $1/4$. Thus the active block certificate is $m=1$, $\kappa=17/20$, and constant envelope. All finite-prefix centers through depth $h$ in both modes form a global Borel cover of cardinality $2^{h+2}-2$ and radius $\sqrt{1800}\ell_h$. The gain estimates follow by comparing the first distinguishing level: its original separation lies between $\ell_{k-1}/6$ and $\ell_{k-1}$, while deletion or prepending shifts this level. These facts verify the hypotheses independently of any claimed memory bound.

The companion proves the resulting exact finite-acquisition law, including the physical-oracle lower and the mode-dependent affine terminal Bernoulli task. If $H_T$ is the actual acquired depth, then the actual score law is a \emph{finite atomic} mixture of the fair-prefix means. Conditional tail variance is comparable to $\E\ell_{H_T}^2$, and a physical Cantor small-ball witness gives the memory term $\ell_{\lfloor\log_2M\rfloor}^2$. Their maximum and sum are both comparable to
\[
 \E\ell_{\min(H_T,\lfloor\log_2M\rfloor)}^2.
\]
The matched sign-minimax law adds $\delta^2$ and charges $B+T+2$ raw calls. The physical law is used solely as an oracle lower witness, never substituted for the actual atomic acquired law. Its proof and all additional assumptions are retained in full in the companion. The progressive acquisition theorem here is a different result: it permits adaptive allocation across coordinates and even smooth factors, but does not replace this nonreset dynamical theorem.

For completeness, exact independent holds preserve the active error radius without requiring a uniform positive activation probability. If active updates satisfy $e_+\le\sqrt\kappa e+\Delta$ and holds copy exactly with independent probability $1-\vartheta_t$, then
\[
 e_t^2\le(1-\vartheta_t)e_{t-1}^2+
       \vartheta_t(\sqrt\kappa e_{t-1}+\Delta)^2.
\]
Both terms preserve $e\le\max\{e_0,\Delta/(1-\sqrt\kappa)\}$. The envelope argument is identical with indicator-form expectations. No conditioning on a null active event is needed when $\vartheta_t=0$. State-dependent missingness in Section~\ref{sec:filters} is not such an independent hold.

\subsection{Finite interfaces and the location of cuts}
A publicly chosen terminal index after an encoding cut means that each retained label selects a vector, or a response function, of possible readouts. It does not mean that the encoder saw the selected index before that cut. This is why the continuation function in Example~\ref{ex:query} has dimension $d$ while its terminal answer has only two values. When a clock or an interface tag can be queried after the cut, its information is part of the retained tuple unless the theorem explicitly conditions on that public interface.

An exact continuation challenge can force disjoint state supports. For a tag $J$ with probabilities $p_j$, unreissued after the cut, exact recovery forces at least one disjoint label set of size $m_j$ per tag. If the score's conditional law is $\nu_j$, the precise distortion is
\[
 \inf_{m_j\ge1,\ \sum_jm_j\le K}\sum_jp_jq_{m_j}(\nu_j).
\]
The proof fixes independent coins, observes that one label cannot decode two distinct tags exactly, and applies conditional nearest-center quantization within each support. Disjoint codebooks give the reverse inequality. This inherited tag mechanism is different from the noisy continuation-function bound of Lemma~\ref{lem:cut}; neither licenses an unchallenged simulator-state lower bound.

On the analytic side all report versions used by wrong candidates must be jointly Borel on the declared true/candidate domain. In the Gaussian case emission densities are positive and the update is everywhere defined on the open simplex. In the continuous-instrument proposition only the positive-density domain is normalized. The root allowance $u$ in Theorem~\ref{thm:transfer} is valid only for routines still satisfying the relation, not for arbitrary floating-point perturbations of the same size. A calibrated-input perturbation can be included in $u$ only after proving that same stable-metric and relation certificate.

These conventions also delimit the retained companion statements. Finite-dimensional identities there do not imply unbounded-operator closure; finite stopped total variation does not imply a large-deviation principle; and same-controller coupling does not imply unrestricted optimal control. The independent realization articles on ordered measurements, finite physical actions and streaming space are not premises of the foundation theorem.

\begin{thebibliography}{99}
\bibitem{anjarlekar} A. Anjarlekar, S. R. Etesami and R. Srikant, \emph{Scalable policy-based RL algorithms for POMDPs}, Advances in Neural Information Processing Systems \textbf{38} (2025).
\bibitem{inductive} R. Andriushchenko, M. \v{C}e\v{s}ka, S. Junges and J.-P. Katoen, \emph{Inductive synthesis of finite-state controllers for POMDPs}, Proc. UAI, Proceedings of Machine Learning Research \textbf{180} (2022), 85--95.
\bibitem{blackwell} D. Blackwell, \emph{Equivalent comparisons of experiments}, Ann. Math. Statist. \textbf{24} (1953), 265--272.
\bibitem{nonanticipative} C. D. Charalambous, P. A. Stavrou and N. U. Ahmed, \emph{Nonanticipative rate distortion function and relations to filtering theory}, IEEE Trans. Automat. Control \textbf{59} (2014), 937--952.
\bibitem{diaconis} P. Diaconis and D. Freedman, \emph{Iterated random functions}, SIAM Rev. \textbf{41} (1999), 45--76.
\bibitem{equitz} W. H. R. Equitz and T. M. Cover, \emph{Successive refinement of information}, IEEE Trans. Inform. Theory \textbf{37} (1991), 269--275.
\bibitem{fox} B. Fox, \emph{Discrete optimization via marginal analysis}, Management Sci. \textbf{13} (1966), 210--216.
\bibitem{fritz} T. Fritz, \emph{A synthetic approach to Markov kernels, conditional independence and theorems on sufficient statistics}, Adv. Math. \textbf{370} (2020), 107239.
\bibitem{ghomi} M. Ghomi, T. Linder and S. Y\"uksel, \emph{Zero-delay lossy coding of linear vector Markov sources: optimality of stationary codes and near optimality of finite memory codes}, IEEE Trans. Inform. Theory \textbf{68} (2022), 3474--3488.
\bibitem{graf} S. Graf and H. Luschgy, \emph{Foundations of Quantization for Probability Distributions}, Lecture Notes in Mathematics 1730, Springer, 2000.
\bibitem{grover} D. Grover and C. Dimitrakakis, \emph{Adaptive belief discretization for POMDP planning}, arXiv:2104.07276, 2021.
\bibitem{hudak} D. Hud\'ak, M. F. L. Galesloot, M. Tappler, M. Kure\v{c}ka, N. Jansen and M. \v{C}e\v{s}ka, \emph{Finite-state controllers for (hidden-model) POMDPs using deep reinforcement learning}, arXiv:2602.08734, 2026 (accepted AAMAS 2026 author version).
\bibitem{kara} A. D. Kara and S. Y\"uksel, \emph{Near optimality of finite memory feedback policies in partially observed Markov decision processes}, J. Mach. Learn. Res. \textbf{23} (2022), no.~11, 1--46.
\bibitem{lecam} L. Le Cam, \emph{Asymptotic Methods in Statistical Decision Theory}, Springer, New York, 1986.
\bibitem{psr} M. L. Littman, R. S. Sutton and S. Singh, \emph{Predictive representations of state}, Advances in Neural Information Processing Systems \textbf{14} (2001), 1555--1561.
\bibitem{nayak} A. Nayak, \emph{Optimal lower bounds for quantum automata and random access codes}, Proc. 40th IEEE Symposium on Foundations of Computer Science (1999), 369--376.
\bibitem{orlitsky} A. Orlitsky and J. R. Roche, \emph{Coding for computing}, IEEE Trans. Inform. Theory \textbf{47} (2001), 903--917; conference version, Proc. FOCS (1995), 502--511.
\bibitem{pages} G. Pag\`es and H. Pham, \emph{Optimal quantization methods for nonlinear filtering with discrete-time observations}, Bernoulli \textbf{11} (2005), 893--932.
\bibitem{companion} Q. Qi, \emph{Supplement S to General Theta Foundations I: Acquired Geometry and Causal Resource Transfer}, submitted supplement to this article, 2026, Sections~1--7; supplied in full as part of the same submission, not a separate publication.
\bibitem{sellami} A. Sellami, \emph{Quantization based filtering method using first order approximation}, SIAM J. Numer. Anal. \textbf{47} (2010), 4711--4734.
\bibitem{shalizi} C. R. Shalizi and J. P. Crutchfield, \emph{Computational mechanics: pattern and prediction, structure and simplicity}, J. Stat. Phys. \textbf{104} (2001), 817--879.
\bibitem{ais} J. Subramanian, A. Sinha, R. Seraj and A. Mahajan, \emph{Approximate information state for approximate planning and reinforcement learning in partially observed systems}, J. Mach. Learn. Res. \textbf{23} (2022), no.~12, 1--83.
\bibitem{torgersen} E. Torgersen, \emph{Comparison of Statistical Experiments}, Cambridge University Press, 1991.
\bibitem{watanabe} S. Watanabe, S. Kuzuoka and V. Y. F. Tan, \emph{Non-asymptotic and second-order achievability bounds for coding with side-information}, IEEE Trans. Inform. Theory \textbf{61} (2015), 1574--1605; accepted author version arXiv:1301.6467v5.
\bibitem{wood} R. G. Wood, T. Linder and S. Y\"uksel, \emph{Optimal zero delay coding of Markov sources: stationary and finite memory codes}, IEEE Trans. Inform. Theory \textbf{63} (2017), 5968--5980.
\bibitem{wynerziv} A. D. Wyner and J. Ziv, \emph{The rate-distortion function for source coding with side information at the decoder}, IEEE Trans. Inform. Theory \textbf{22} (1976), 1--10.
\bibitem{zhang2012} Z. Zhang, M. Littman and X. Chen, \emph{Covering number as a complexity measure for POMDP planning and learning}, Proc. AAAI \textbf{26} (2012), 1853--1859.
\bibitem{zhang2014} Z. Zhang, D. Hsu and W. S. Lee, \emph{Covering number for efficient heuristic-based POMDP planning}, Proc. ICML, Proceedings of Machine Learning Research \textbf{32} (2014), 28--36.
\bibitem{zhu-moran} S. Zhu, \emph{On the asymptotic quantization error for the doubling measures on Moran sets}, arXiv:1908.00202, 2019.
\bibitem{zhulu} Y. Zhu and Y. Lu, \emph{A covering framework for offline POMDPs learning using belief space metric}, Proc. AISTATS, Proceedings of Machine Learning Research \textbf{300} (2026), 1423--1431.
\end{thebibliography}

\end{document}
